% concatenated from R-triviality.tex
\documentclass[12pt]{amsart}
\usepackage{amssymb,amsmath,amsfonts,latexsym}
\usepackage{mathtools}
\usepackage{tikz-cd}
\usepackage[colorlinks=true, linkcolor=blue, citecolor=red, anchorcolor=orange,]{hyperref}
%\usepackage{bm, enumerate}
% \usepackage{xcolor}

\setlength{\textheight}{9in}\setlength{\textwidth}{475pt}
\oddsidemargin -0mm \evensidemargin -0mm \topmargin -0pt
\newcommand{\newsection}[1]{\setcounter{equation}{0} \section{#1}}
\setcounter{footnote}{1}
\renewcommand{\theequation}{\arabic{section}.\arabic{equation}}
%%%%%%%%%%%%%%%%%%%%%%%%%%%Shortenings%%%%%%%%%%%%%%%%%%%

\newcommand{\vp}{\varphi}
\newcommand{\cla}{\mathcal{A}}
\newcommand{\clb}{\mathcal{B}}
\newcommand{\clc}{\mathcal{C}}
\newcommand{\cld}{\mathcal{D}}
\newcommand{\cle}{\mathcal{E}}
\newcommand{\clf}{\mathcal{F}}
\newcommand{\clg}{\mathcal{G}}
\newcommand{\clh}{\mathcal{H}}
\newcommand{\cli}{\mathcal{I}}
\newcommand{\clj}{\mathcal{J}}
\newcommand{\clk}{\mathcal{K}}
\newcommand{\cll}{\mathcal{L}}
\newcommand{\clm}{\mathcal{M}}
\newcommand{\cln}{\mathcal{N}}
\newcommand{\clo}{\mathcal{O}}
\newcommand{\clp}{\mathcal{P}}
\newcommand{\clq}{\mathcal{Q}}
\newcommand{\clr}{\mathcal{R}}
\newcommand{\cls}{\mathcal{S}}
\newcommand{\clt}{\mathcal{T}}
\newcommand{\clu}{\mathcal{U}}
\newcommand{\clw}{\mathcal{W}}
\newcommand{\clx}{\mathcal{X}}
\newcommand{\cly}{\mathcal{Y}}
\newcommand{\clz}{\mathcal{Z}}


\newcommand\Zp{\mathbb{Z}_+^n}
\newcommand{\D}{\mathbb{D}}
\newcommand{\N}{\mathbb{N}}
\newcommand{\C}{\mathbb{C}}
\newcommand{\T}{\mathbb{T}}
\newcommand{\Z}{\mathbb{Z}}
\newcommand{\z}{\bm{z}}
\newcommand{\w}{\bm{w}}
\newcommand{\raro}{\rightarrow}
\newcommand{\Raro}{\Rightarrow}
\newcommand{\subnq}{\begin{array}[t]{c} \subset \\*[-15pt] \scriptstyle\neq
\end{array}}
\newcommand{\supnq}{\begin{array}[t]{c} \supset \\*[-15pt] \scriptstyle\neq
\end{array}}
\newcommand{\vsp}{\vspace{5mm}}
\newcommand{\hsp}{\hspace{5mm}}
\newcommand{\vspp}{\vspace{1cm}}
\def \qed {\hfill \vrule height6pt width 6pt depth 0pt}
\def\textmatrix#1&#2\\#3&#4\\{\bigl({#1 \atop #3}\ {#2 \atop #4}\bigr)}
\def\dispmatrix#1&#2\\#3&#4\\{\left({#1 \atop #3}\ {#2 \atop #4}\right)}
\newcommand{\be}{\begin{equation}}
	\newcommand{\ee}{\end{equation}}
\newcommand{\ben}{\begin{eqnarray*}}
	\newcommand{\een}{\end{eqnarray*}}
\newcommand{\CC}{\centerline}
\newcommand{\NI}{\noindent}
\newcommand{\addeqn}{\addtocounter{equation}{1}}
\def \pf {{\sc Proof}\hspace{.5cm}}
\newcommand{\bi}{\begin{itemize}}
	\newcommand{\ei}{\end{itemize}}
\newcommand{\diag}{\mbox{diag}}
\def\npair{{\mbox{\boldmath $n$}}}

\newcommand\la{{\langle }}
\newcommand\ra{{\rangle}}






%%%%%%%%%%%%Theorem Environment%%%%%%%%%%%%%%%%%%%%%%
\newtheorem{thmI}{Theorem}
\renewcommand{\thethmI}{\Alph{thmI}}
\newtheorem{theorem}{Theorem}[section]
\newtheorem{lemma}[theorem]{Lemma}
\newtheorem{proposition}[theorem]{Proposition}
\newtheorem{corollary}[theorem]{Corollary}
\theoremstyle{definition}
\newtheorem{definition}[theorem]{Definition}
\newtheorem{example}[theorem]{Example}
\newtheorem{exercise}[theorem]{Exercise}
\newtheorem{conclusion}[theorem]{Conclusion}
\newtheorem{conjecture}[theorem]{Conjecture}
\newtheorem{criterion}[theorem]{Criterion}
\newtheorem{summary}[theorem]{Summary}
\newtheorem{axiom}[theorem]{Axiom}
\newtheorem{question}[theorem]{Question}
%\theoremstyle{remark}
\newtheorem{remark}[theorem]{Remark}
\numberwithin{equation}{section}

\begin{document}
	
	
\title{On $R$-equivalence of Automorphism groups of Associative algebras}
	
	
	
\author[Das]{Dibyendu Das}
\address{Indian Statistical Institute, Statistics and Mathematics Unit, 8th Mile, Mysore Road, Bangalore, 560059,
India}
\email{ddas4282@gmail.com}
	
	
%\today
	
	
\subjclass{20G15, 14E08, 16Gxx}
\keywords{Automorphism groups of algebras, Associative algebras, Rationality, Stably rational, $R$-triviality, Quiver and relations.}
	
\begin{abstract}
Let $A$ be a finite-dimensional associative $k$-algebra with identity. The primary aim of this paper is to study the rationality properties of the group of all $k$-algebra automorphisms of $A$, as an affine algebraic group over an arbitrary field $k$. We investigate mainly the $R$-equivalence property of the identity component of $\mathrm{Aut}_{k}(A)$ over a perfect field $k$.
\end{abstract}
	
\maketitle


\tableofcontents

	
\newsection{Introduction and Main results }\label{sec: intro}

Let $k$ be an arbitrary field and $A$ be a finite-dimensional unital associative algebra over $k$. We denote by $\text{Aut}_k(A)$ the group of all $k$-algebra automorphisms of $A$.
The understanding of the algebraic group $\text{Aut}_k(A)$ and its different properties is of deep interest. Many researchers have looked into the group of automorphisms of finite-dimensional unital associative algebras from an algebraic perspective. For example, see the papers (\cite{Djokovic}, \cite{Luks}, \cite{H.G}, \cite{Stanley}, \cite{MM}) and their references for more details. The papers \cite{Frolich} and \cite{Yau2} discuss relations of $\text{Aut}_k(A)$ with the Picard group of $A$ and the classification of isolated hypersurface singularities, respectively. The article (\cite{Pollack}) by David R. Pollack served as our main inspiration, where he has shown how different properties of $G_A$ (the identity component of $\text{Aut}_k(A)$) affect the properties of $A$ over an algebraically closed field. The solvability and unipotency of the group $\text{Aut}_k(A)$ are important in this context, as is apparent from Yau's work in (\cite{Yau2}). Perepechko (\cite{Perepechko}) and Saor\'in et al. (\cite{AS1}, \cite{AS2}) have made significant contributions to the understanding of the solvability of $\text{Aut}_k(A)$. Moreover, Saor\'in and Guil Asensio explored the Picard groups along with the outer automorphism groups of finite-dimensional associative algebras in their papers (\cite{AS3}, \cite{AS2}). Nonetheless, little is known about the geometric properties of $\text{Aut}_k(A)$. This paper studies some birational properties of $\text{Aut}_k(A)$ as a $k$-group. One such property is similar to the concept of path connectedness. In algebraic geometry, this notion corresponds to the $R$-equivalence of a variety. We have explored the rationality and $R$-equivalence properties of $G_A$ (the identity component of $\mathrm{Aut}_k(A)$), which is crucial in the study of $\text{Aut}_k(A)$.

\vspace{0.2in}
\noindent
For a connected algebraic group $G$ defined over an arbitrary field $k$, an important question is the rationality problem: Is the function field $k(G)$ purely transcendental over $k$? Determining whether a variety is rational is a central problem in algebraic geometry. In this context, the notion of $R$-equivalence was introduced by Y.~Manin around $1972$ (see \cite{Manin}, Chapter $2$). The knowledge that a group is $R$-trivial helps in studying its rationality properties. A systematic study of this notion for linear algebraic groups was then initiated by J.-L.~Colliot-Thélène and J.-J.~Sansuc \cite{CTS}. Around 1996, A.~Merkurjev \cite{AM2} and Ph.~Gille \cite{Gille2} studied $R$-triviality for the group of proper projective similitudes of a quadratic form or, more generally, automorphism groups of central simple associative algebras with involutions. For more on $R$-equivalence of automorphism groups of central simple algebras with involutions, one can look at \cite{PA}, \cite{AP}, \cite{Preeti} and their references. On the other hand, the automorphism groups of finite-dimensional associative algebras are not as thoroughly investigated in this context.
\vskip2mm 
\begin{question}
Let $G=G_A$, where $A$ is a finite-dimensional associative algebra over a perfect field $k$. Is the function field $k(G_A)$ purely transcendental over $k$?  
\end{question}
\vskip2mm
To the best of our knowledge, there are no known results regarding the above question. $R$-triviality is the necessary condition for a group to be rational (see Section \ref{R-equivalence}). This article provides a systematic framework for investigating the $R$-equivalence property of connected algebraic groups that are of the form of $G_A$, where $A$ is a split finite-dimensional associative algebra, which is essential for the investigation of any finite-dimensional associative algebra (see Section \ref{sec: Preliminaries} and Definition \ref{Definition 2.6}). For simplicity, our results are given for the case of split finite-dimensional associative algebras. More precisely, we have the following results for $G_A$ over a perfect field $k$, where $A$ is a split finite-dimensional associative algebra:
\vspace{0.2in}

\begin{thmI}[]
$(=\text{Theorem }\ref{Theorem 4.20})$
Let $A$ be a split finite-dimensional associative algebra over a perfect field $k$ with $A_s$ a semisimple subalgebra of $A$ such that $A= A_s\oplus J$ (vector space decomposition) and $G_{A, A_s}\coloneqq\{\sigma\in G_A: \sigma(a)=a \quad \forall a\in A_s\}$. Then, $G_A$ is $R$-trivial if and only if $G_{A,A_s}$ is $R$-trivial. Suppose any one of the following holds:
\begin{enumerate}
    \item $J^2=0$;
    \vskip2mm
    \item $\emph{dim}(J/J^2)\leq 5$;
    \vskip2mm
    \item $\emph{rank} (G_{A,A_s})=\emph{dim}(J/J^2)$.
\end{enumerate}
\vskip2mm
 Then, $\mathrm{Aut}_k(A)^{0}\coloneqq G_A$ is $R$-trivial.
\end{thmI}
\vskip2mm

\begin{thmI}[]
$(=\text{Theorem }\ref{Theorem 4.33})$ Let $A$ be a split finite-dimensional associative algebra over a perfect field $k$. Assume that $\emph{\text{Aut}}_k(A)$ is a nilpotent group. Then, $\emph{\text{Aut}}_k(A)^{0}\coloneqq  G_A$ is rational.
\end{thmI}

\vspace{0.2in}
\noindent
At first, we reduce the problem of studying $R$-equivalence of $G_A$, where $A$ is a split associative algebra, to that of studying $R$-equivalence of $G_A$, where $A$ is a split local algebra (see Definition \ref{Definition 2.6}).
\vskip2mm 
\noindent
%\textbf{Reduction to local algebras:}
We prove in Theorem \ref{Theorem 4.20} that $G_A$ is $R$-trivial if and only if $G_{A,A_s}$ is $R$-trivial, where $A_s$ is a semisimple subalgebra of $A$ such that $A=A_s\oplus J$ (vector space decomposition) with Jacobson radical $J$ and $G_{A,A_s}\coloneqq\{\sigma\in G_A: \sigma(a)=a \quad \forall a\in A_s\}$. In Proposition \ref{Proposition 4.34}, we show that the property of $R$-triviality of $G_A$ is invariant under \emph{Morita equivalence} (see Definition \ref{Definition 2.3}) using Theorem \ref{Theorem 4.20}. Consequently, it suffices to examine split local algebras in order to investigate $R$-triviality of $G_A$ as shown in Remark \ref{Remark 4.35}. Furthermore, Proposition \ref{Proposition 4.37} indicates that for a split local algebra $A$, the group $G_A$ is stably birationally equivalent to $\text{ Im}(\Phi_A)$, where $\Phi_A: G_A\rightarrow \mathrm{GL}(J/J^2)$ is the canonical map and $J$ is the Jacobson radical of $A$. In particular, we have \[G_A(F)/R\cong \mathrm{Im}(\Phi_A)(F)/R\] for every field extension $F/k$ and for every split local algebra $A$ over $k$ (see Section \ref{R-equivalence} for notation). Therefore, the analysis can be effectively narrowed down to the study of $\text{ Im}(\Phi_A)$.
\vskip4mm
\noindent
Thereafter, we are primarily concerned with the study of $R$-triviality of the automorphism groups of split local commutative algebras. In Lemma \ref{Lemma 5.33}, we show the existence of a $\mathrm{Im}(\Phi_A)$-stable finite-dimensional subspace $W$ associated to a split local commutative algebra $A$ using \emph{quiver and relation} representation of $A$. Then we use this subspace to examine the $R$-equivalence property of $G_A$, where $A$ is a split local commutative algebra. Under certain new sufficient conditions, we prove that $G_A$ is $R$-trivial, thereby establishing that the conditions of Theorem \ref{Theorem 4.20} are sufficient but not necessary for commutative algebras (see Example \ref{Example 5.15} and Theorem \ref{Theorem 5.35}). Furthermore, we show the existence of a local commutative algebra $A$ with $\text{dim}(J/J^2)\geq 6$ and $G_A$ is not $R$-trivial (see Theorem \ref{Theorem 6.4}). More precisely, we have the following results stated below for $G_A$ over a perfect field $k$, where $A$ is a split local commutative algebra:
\vspace{0.2in}

\begin{thmI}
$(=\text{Theorem }\ref{Theorem 5.13})$ Let $A$ be a finite-dimensional local commutative algebra over a field $k$ with $\text{dim}(J/J^2)=n$. Suppose $A\cong k[X_1,\dots, X_n]/I$ by quiver and relation representation of $A$, where $I=\langle X_1,\dots, X_n\rangle^l+\langle f\rangle$ and $2\leq \mathrm{ deg}(f)<l$ is the Lowey length of $A$. Then, $G_A$ is rational or $R$-trivial as a $k$-group if and only if $\mathrm{Stab}^{0}(f)$ is rational or $R$-trivial.
\end{thmI}
\vskip2mm
\begin{thmI}
$(=\text{Theorem }\ref{Theorem 5.35})$ Let $A$ be a split local commutative graded by the radical algebra over a field $k$ with $\emph{dim}(J/J^2)=n$. Let $W$ be the minimal degree subspace associated to $A$ (see Remark \ref{Remark 5.34}). Assume $W$ contains a non-singular homogeneous polynomial of degree at least $3$ if $\emph{char }k=0$. Additionally if $\emph{char }k=p>3$ and $W$ contains a non-singular homogeneous polynomial of degree $d$ with $2<d<p$. Then the following two statements are equivalent:
\vskip2mm
\begin{enumerate}
\item $G_A$ is $k$-split solvable; 
\vskip2mm
\item $W$ has a stable full flag by $\emph{\text{Im}}(\Phi_A)$.
\end{enumerate}
Therefore, under the condition $(2)$, the group is rational. In particular, if the $\text{dim }W=1$, then $G_A$ is rational.
\end{thmI}
\vskip2mm
\begin{thmI}
$(=\text{Theorem }\ref{Theorem 6.4})$ For each $n\geq3$, there exists a finite-dimensional commutative algebra $A_n$ over a field $k_n$ with $\emph{dim } \displaystyle (\frac{J_n}{J_n^2})\geq 2n$, where $J_n$ is the Jacobson radical of $A_n$, such that $G_{A_n}$ is not $R$-trivial as a $k_n$-group and rank of $G_{A_n}$ is $(n+1)$. 
\end{thmI}
\vspace{0.2in}

Apart from the above geometric property of $\text{Aut}_k(A)$, we have also explored some algebraic properties of $\text{Aut}_k(A)$. In Theorem \ref{Theorem 3.4}, we show that if $\text{Aut}_k(A)\cong T$, where $T$ is a torus defined over a perfect field $k$ and $A$ is a split finite-dimensional associative algebra over $k$, then $T$ must be split. One knows that split groups are always rational. Towards rationality, it is of interest to know if the identity component of the automorphism group $G_A$ is split when $A$ is a split local associative algebra. In Proposition \ref{Proposition 4.12}, we prove that for a split local associative algebra $A$, the rank of $G_A$ is bounded by $\text{dim}(J/J^2)$.  From the results of \cite{AS1}, it follows that $G_A$ is split for a split local commutative monomial algebra $A$ with rank of $G_A$ equals to $\text{dim}(J/J^2)$. In \cite{AS2}, these results have been refined for semimonomial algebras. Nevertheless, this assertion does not hold in all cases. We present a set of examples where the group is not split (see Theorem \ref{Theorem 5.16} and Corollary \ref{Corollary 5.19}). Saor\'in et al. have previously shown in \cite{AS1} that if $A$ is a split local commutative algebra with $\text{ dim}(J/J^2)=n$, then $G_A$ contains a split torus of rank $n$ if and only if $A$ is a monomial algebra. This motivates us to explore a generalisation of this question. Namely:
\vskip2mm
\begin{question}
Let $A$ be a split local commutative algebra over field $k$ with $\text{dim}(J/J^2)=n$. What is the necessary and sufficient condition on $A$ so that $G_A$ contains a split torus of rank $r< n=\text{dim}(J/J^2)$? 
\end{question}
\vskip2mm
\noindent
We answer this question in Theorem \ref{Theorem 5.25}. The more precise results are given below:
\vspace{0.2in}

\begin{thmI}
$(=\text{Theorem }\ref{Theorem 3.4})$ Let $T$ be a torus of rank $r$ over a perfect field $k$ such that there exists a split finite-dimensional associative algebra $A$ over $k$ with $\emph{\text{Aut}}_k(A)\cong T$, then $T$ is split.   
\end{thmI}
\vskip2mm
\begin{thmI}
$(=\text{Theorem }\ref{Theorem 5.16})$ Let $A$ be a finite-dimensional local commutative algebra with $\text{dim}(J/J^2)=n$. Suppose $A\cong k[X_1,\dots, X_n]/I$ by quiver and relation representation of $A$, where $I$ is given by $\langle X_1, X_2,\dots, X_n\rangle^l+\langle q\rangle$ and $\emph{char }k\neq2$; $l>2$ is the Lowey length of $A$ and $(J/J^2,q)$ is an anisotropic quadratic space over the field $k$, then $G_A$ is not a $k$-split group.
\end{thmI}
\vskip2mm
\begin{thmI}
$(=\text{Theorem } \ref{Theorem 5.25})$ Let $A$ be a split local commutative algebra over an infinite field $k$ with $\emph{dim}(J/J^2)=n$. Let $r\in \mathbb{N}$ and $1\leq r\leq n$. Then the following statements are equivalent:
\begin{enumerate}
    \vskip4mm
    \item $G_A$ contains a diagonal change of variables of the form \[D(r,k)=
\begin{pmatrix}
a_1 & 0 &\cdots & 0 & 0 & \cdots & 0 \\
0 & a_2 & \cdots & 0 & 0 & \cdots & 0\\
\vdots & \vdots & \ddots & \vdots& \vdots & \vdots & \vdots\\
0 & 0 & \cdots & a_{r-1} & 0 &\cdots & 0 \\
0 & \cdots & 0 & 0 & \beta & \cdots & 0\\
\vdots & \vdots & \vdots & \vdots & \vdots & \ddots & \vdots\\
0 & \cdots & \cdots & 0 & 0 & \cdots & \beta
\end{pmatrix}_{n\times n}
,\]
where  $a_{i}, \beta \in k^{\times}, 1\leq i\leq r-1;$
\vskip4mm
     \item There exists an admissible ideal $I$ which satisfies the Property $*$ for $r$ (see Definition \ref{Definition 5.24}) or a monomial ideal $I$ such that $A\cong k[X_1,\dots,X_n]/I$.
\end{enumerate}
\vskip4mm
If the hypothesis $(2)$ is satisfied, we have a $k$-embedding $\phi: \mathbb{G}_m^{r}\rightarrow G_A$. 
\end{thmI}
\vskip2mm
\noindent
Clearly $\text{Aut}_{k}(A)$ is an affine algebraic group. Associative algebras and affine algebraic groups have a rich structure theory. From now onwards, algebra means a finite-dimensional unital associative algebra. Both (\cite{Humphreys}) and (\cite{Springer}) serve as general references for algebraic groups, and we use the theory of associative algebras from (\cite{Pierce}, \cite{Auslander}, \cite{Curtis}). The concept of representation of a split basic finite-dimensional associative algebra by \emph{quiver and relation} is taken from (\cite{Auslander}, \cite{PG}, \cite{ISD}). For the theory of Galois cohomology and quadratic forms, one can follow \cite{KMRT}. In this article, we will deal with finite-dimensional associative algebras over perfect fields (some of the results may be valid over arbitrary fields, but we mainly stick to perfect fields). Unless specified otherwise, there are no other restrictions on the field $k$.
\vskip2mm

The rest of the paper is organized as follows: In Section \ref{sec: Preliminaries}, we recall some well-known results on the structure of finite-dimensional associative algebras and provide some basic definitions that are used throughout the paper. Section \ref{Review} reviews Pollack's results over perfect fields. Section \ref{R-equivalence} is devoted to the study of $R$-equivalence of the automorphism groups of split finite-dimensional associative algebras. In Section \ref{sec: commutative algebra}, we establish the result about the $R$-equivalence of automorphism groups of split local commutative algebras by using quiver and relation. Finally, in Section \ref{counter}, we provide examples that demonstrate the existence of non-$R$-trivial groups of the form $G_A$, where $A$ is a finite-dimensional commutative algebra.
 
\section{Preliminaries}\label{sec: Preliminaries}
\noindent
This section discusses the structure theory of finite-dimensional associative algebras and provides some basic definitions that are necessary for subsequent sections.
\vskip2mm
\noindent
Throughout this paper, $A$ is a finite-dimensional associative algebra over a field $k$ with unity, and $J$ is its Jacobson radical (in short, we call it radical). For any affine algebraic group $G$ over $k$, $G^{0}$ denotes its identity component. In this article, $\mathrm{Aut}_k(A)^{0}=G_A$. Following the point of view in (\cite{KMRT}), we can regard any affine algebraic group as a representable functor from the category of algebra to the category of groups. In this sense $\text{Aut}_k(A)$ is a functor, 
\begin{equation*}\label{eq1}
\text{Aut}_k(A): B\mapsto \text{Aut}_{B}(A\otimes_{k} B)
\end{equation*}
Let $G$ be an algebraic group over $k$, then $G(F)$ denotes its $F$-rational points, where $F$ is a field extension of $k$.
\vskip2mm
\noindent
Next, we recall the principal structure theorem for finite-dimensional associative algebras: the Wedderburn-Malcev decomposition.

\vspace{2mm}

\begin{theorem}\label{Theorem 2.1}
\textbf{$($\emph{Wedderburn-Malcev}}, \cite{Curtis}, \emph{Theorem} $72.19)$ Let $A$ be a finite-dimensional associative algebra over $k$ (arbitrary field) with radical $J$ such that the residue class algebra $B=A/J$ over $k$ is separable. Then, there exists a semisimple subalgebra $A_s$ of $A$ such that $A=A_s\oplus J$ (vector space decomposition). If there exists a semisimple subalgebra $A'_s$ such that $A=A'_s\oplus J$, then there exists an element $r\in J$ such that $A'_s=(1+r)A_s(1+r)^{-1}$.
\end{theorem}

\vskip2mm

\begin{corollary}\label{Corollary 2.2}
$($\cite{Pierce}, \emph{Page} $211)$ If $k$ is a perfect field, and $A$ is finite-dimensional associative $k$-algebra, then there is a semisimple subalgebra $A_s$ of $A$ such that $A/J\cong A_s$ and $A=A_s\oplus J$. Moreover, $A_s$ is unique upto
conjugation by units of the form $(1+r)$, where $r\in J$.
\end{corollary}

\vskip2mm 
\noindent
Let $A^{\times}$ denote the group of units in $A$ and $k^{\times}=k-\{0\}$. We have an obvious map from $A^{\times}$ to $\text{Aut}_k(A)$ given  by $\psi:A^{\times}\rightarrow \text{Aut}_k(A) \text{ with }\psi(a)=\text{Int}(a)$, where $\text{Int}(a)$ is the automorphism  $x\mapsto axa^{-1}$.

\vskip2mm
\noindent
This is a morphism of algebraic groups, whose image is $\text{Inn}(A$), the group of all inner automorphisms of $A$. Observe that $\{1+r:r\in J\}$ is a closed, connected, unipotent subgroup of $A^{\times}$ whose image $\hat{J}\coloneq\{\text{Int}(1+r):r\in J\}$ is normal in $\text{Aut}_k(A)$ and $\hat{J}\subset U$, where $U$ is the $k$-unipotent radical of $\text{Aut}_k(A)$.

\vskip2mm

\begin{definition}\label{Definition 2.3}$($\cite{Pierce}, Section $9.6$$)$
Two algebras $A$ and $B$ are said to be \emph{Morita equivalent} if their left module categories are equivalent.
\end{definition}

\vskip2mm
\begin{definition}\label{Definition 2.4}
An algebra $A$ over $k$ is a \emph{basic algebra} if $A/J\cong \bigoplus_{i=1}^{n} D_i$, where each $D_i$ is a division algebra over $k$ and $n\in \mathbb{N}$. In particular, if $A/J\cong D$ for some division algebra $D$ over $k$, we call $A$ a \emph{local algebra} over $k$.  
\end{definition}

\vskip2mm 
\begin{theorem}\label{Theorem 2.5}
$($\cite{Pierce}, \emph{Section} $9.6)$ For each finite-dimensional associative algebra $A$ over a perfect field $k$, there always exists a basic algebra $B$ of $A$ which is Morita equivalent to $A$.   
\end{theorem}
\vskip2mm
\noindent
Thus, over a perfect field $k$, every finite-dimensional associative algebra \( A \) can be decomposed as \( A = A_s \oplus J \) (vector space decomposition), as established by Corollary \ref{Corollary 2.2}. In this decomposition, \( A_s\cong A/J \) is a semisimple subalgebra that is isomorphic to \( \bigoplus_{i=1}^{m} M_{n_{i}}(D_{i}) \), where each $D_i$ is a division algebra over $k$ for all $1\leq i\leq m$ and \( J \) denotes the Jacobson radical of the algebra \( A \). 
\vskip2mm
\begin{definition}\label{Definition 2.6}
An associative algebra $A$ over $k$ is called a \emph{split associative algebra} if $A/J\cong \bigoplus_{i=1}^{m} M_{n_{i}}(k)$, where $m\in \mathbb{N}$. An associative algebra $A$ over $k$ is called a \emph{split basic associative algebra} if $A/J\cong \bigoplus_{i=1}^{m}k$, and an associative algebra $A$ over a field $k$ is a \emph{split local associative algebra} if $A/J\cong k$.  
\end{definition}

\vskip2mm
\section{Review of Pollack's Results over Perfect fields}\label{Review}
\vskip2mm
\noindent
In this section, we discuss the results of (\cite{Pollack}) related to the reductivity and semisimplicity of $G_A$ in terms of the properties of $A$ over a perfect field $k$, where $A$ is a split associative algebra over $k$. Over an algebraically closed field, any finite-dimensional associative algebra is split. Therefore, for a split finite-dimensional associative algebra over a perfect field $k$, the proof of Theorem \ref{Theorem 3.1} follows from a similar line of argument as done by Pollack in \cite {Pollack}. The other proofs of this section are given for the sake of completeness. Let us recall a few notations, $\text{Inn}(A)$ is the inner automorphism group of $A$, $\hat{J}=\{\text{Int}(1+r):r\in J\}$ and $U$ denotes the $k$-unipotent radical of $G_A$.

\vskip2mm
\begin{theorem}\label{Theorem 3.1}
Let $A$ be a split finite-dimensional associative algebra over a perfect field $k$, and let $J$ be the radical of $A$. If $U$ is the $k$-unipotent radical of $G_A$, then $J^2=0$ if and only if $J\subset A^{U}$, where $A^U\coloneq \{a\in A:\sigma(a)=a \quad \forall \sigma\in U(F), \text{ for all field extension $F$ over $k$ } \}$. 
\end{theorem}
\vskip1mm
\begin{proof}
Since over a perfect field $k$, a group is reductive if and only if the $k$-unipotent radical is trivial, and we also have $A/J\cong \bigoplus_{i=1}^{m}M_{n_{i}}(k)$. Therefore, the proof of the above statement follows in a similar line of argument as done by Pollack in Theorem $1.1$ of \cite{Pollack}.
\end{proof}
\vskip2mm

\vspace{2mm}
\begin{corollary}\label{Corollary 3.2}
Let $A$ be a split finite-dimensional associative algebra over a perfect field $k$. Then $G_A$ is reductive if and only if $J^2=0$ and $J\subset Z(A)$, where $Z(A)$ is the center of $A$.
\end{corollary}
\vskip1mm
\begin{proof}
    Let $G_A$ be reductive, then $U=\{Id\}$, $U$ is the $k$-unipotent radical. Hence, it follows that $A^U=A\supset J$. Therefore, by Theorem \ref{Theorem 3.1}, we have $J^2=0$. Since $U=\{Id\}$, we also have $\hat{J}=\{Id\}$, which further implies $J\subset Z(A)$. Conversely, if we assume $J^2=0$, then it follows from the proof of Theorem \ref{Theorem 3.1} that the $k$-unipotent radical of $G_A$ is $\hat{J}=U$. However, $J\subset Z(A)$ shows that $\hat{J}=\{Id\}$. Hence, the group $G_A$ is reductive.
\end{proof}

\vskip2mm
\begin{corollary}\label{Corollary 3.3}
Let $A$ be a split finite-dimensional associative algebra over a perfect field $k$. Then $G_A$ is semisimple if and only if $A\cong \bigoplus_i M_{n_i}(k)$. 
\end{corollary}
\vskip1mm
\begin{proof}
$G_A$ is semisimple implies that it is reductive, so $J^2=0$ and $J\subset Z(A)$. We will show that $J=0$. Suppose $J\neq 0$, we have a morphism from $G_A$ to $GL(J)$ as follows $\epsilon_A: G_A\rightarrow \text{GL}(J)$, given by $\epsilon_A(\phi)=\phi|_{J}$. Using the fact $J^2=0$, we can prove that the homomorphism is surjective. Let $T\in \text{GL}(J)$, then consider $\Tilde{T}(a+r)=a+T(r)$, where $a\in A_s$ and $r\in J$. However, $J^2=0$ implies $\Tilde{T}\in G_A$ and $\epsilon_A(\Tilde{T})=T$. So, $\epsilon_A$ is a surjection. Therefore, $\text{GL}(J)$ is semisimple. Hence, $J$ must be trivial. Conversely, if $A\cong \bigoplus_i M_{n_i}(k)$, then $G_{A}\cong \prod_i \mathrm{PGL_{n_i}}$, which is a semisimple group.
\end{proof}

\vskip2mm
\begin{theorem}\label{Theorem 3.4}
 Let $T$ be a torus of rank $r$ over a perfect field $k$ such that there exists a split finite-dimensional associative algebra $A$ over $k$ with $T\cong \emph{\text{Aut}}_k(A)$, then $T$ is a split torus.
\end{theorem}
\begin{proof}
It is given that $T\cong \text{Aut}_k(A)=G_A$ over a perfect field $k$, where $A$ is a split finite-dimensional algebra over $k$. Since $T$ is nilpotent, it follows from the proof of Theorem \ref{Theorem 3.1} that $A$ must be a split basic algebra. Therefore, $\prod_i G_{A_i}\cong T$, where $A_i$ is a split local algebra such that $A\cong \bigoplus_i A_i$. Hence, $G_{A_i}$ is either a torus or a trivial group. If $G_{A_i}$ is a trivial group, then $A_i=\bigoplus_j k$, and if $G_{A_i}$ is a torus, then by using Corollary \ref{Corollary 3.2}, we get that $J_i^2=0$ and $J_i\subset Z(A_i)=k\oplus Z(J_i)$, where $J_i$ is the Jacobson radical of $A_i$. Since each $A_i$ is a split local algebra, we have that each $A_i$ is commutative. However, by the quiver and relation representation of $A_i$, we get that $A_i\cong k[X_1,\dots,x_n]/I$ with $\langle X_1,\dots,X_n\rangle^l\subseteq I \subseteq\langle X_1,\dots,X_n\rangle^2$, where $l=\text{min }\{n:J_i^n=0\}\geq 2$ and $n\geq 1$ (see \cite{AS1} or Section \ref{sec: commutative algebra}). Since $J_{i}^2=0$, we have $A_i\cong k[X_1,\dots,X_n]/\langle X_1,\dots, X_n\rangle^2$. Therefore, for each $G_{A_i}$ to be a torus, we must have $A_i\cong k[X]/\langle X\rangle^2$ by Theorem $1.3$ of \cite{AS1}. Finally we have either $A_i\cong \bigoplus_j k$ or $A_i\cong k[X]/\langle X\rangle^2$. Hence, $A\cong (\bigoplus_j k)\oplus (\bigoplus_i k[X]/\langle X\rangle).$ Now it follows that $G_A\cong \mathbb{G}_{m,k}^r$.
\end{proof}

\vskip2mm
\noindent
The following results are related to the radical of an associative algebra $A$ which may not be finite-dimensional; here $rad(A)$ denotes the Jacobson radical of $A$.

\vskip2mm

\begin{theorem}\label{Theorem 3.5}
$($\cite{Lam}, \emph{Theorem} $5.14)$ For any $k$-algebra $A$ and any field extension $K/k$, we have $A\cap rad(A\otimes K)\subset rad A$. If $K/k$ is an algebraic extension, or if $\mathrm{dim}_{k}A < \infty$, then $A \cap rad(A\otimes K)=rad A$.
\end{theorem}

\vskip2mm

\begin{theorem}\label{Theorem 3.6}
$($\cite{Lam}, \emph{Theorem} $5.17)$ Let $A$ be a $k$-algebra and $K/k$ be a separable algebraic extension. Then $rad(A\otimes_{k} K)$ = $(radA)\otimes_{k} K$.
\end{theorem}

\vskip2mm

\begin{remark}\label{Remark 3.7}
   Let $A$ be a split finite-dimensional algebra over the perfect field $k$, with $J$ being the radical of $A$. If $A=A_s\oplus J$, where $A_s\cong \bigoplus_{i=1}^{m} M_{n_i}(k)$, then $A\otimes_k K=(A_s\otimes_k K)\oplus (J\otimes_k K)$, where $A_s\otimes_k K\cong \bigoplus_{i=1}^{m} M_{n_i}(K)$, for every extension $K$ of $k$ and $m\in \mathbb{N}$.
\end{remark}

\vskip2mm
\section{$R$-Equivalence and $R$-triviality}\label{R-equivalence}
\vskip2mm
\noindent
Let $k$ be a perfect field and $G_A$ be the identity component of the automorphism group of a split finite-dimensional associative algebra $A$. In this section, we discuss the $R$-equivalence of $G_A$.

\vskip2mm
\noindent
Let $X$ and $Y$ be varieties (irreducible) over $k$. We write $X\cong_b Y$ if $X$ and $Y$ are birationally equivalent, that is, the function fields $k(X)$ and $k(Y)$ are isomorphic over $k$. And we use the notation $X \cong^{s.b} Y$ if $X$ and $Y$ are stably birationally equivalent, that means, $X\times \mathbb{A}_k^m\cong_b Y\times \mathbb{A}_k^n$ for some $m,n$.

\vskip2mm
\noindent
We say that a variety $X$ is \emph{rational} (resp. \emph{stably rational)} if $X\cong_b\mathbb{A}_k^n$ (resp. $X\times_k \mathbb{A}_k^n\cong_b \mathbb{A}_k^{m})$  for some $n,m\geq 0$.
  
\vskip2mm
\noindent
Let X be an irreducible variety over a field k with $X(k)\neq\phi$. We define two points $x, y \in X(k)$ to be $R$-equivalent ($x \sim y$) if there exists a sequence points $x_0= x, x_1,...,x_n=y$ in $X(k)$ and rational
maps $f_i:\mathbb{A}^{1}_{k}\dashrightarrow X$, $1\leq i\leq n$, defined over $k$, regular at $0$ and $1$, such
that $f_{i}(0)=x_{i-1},f_{i}(1)=x_{i}$ (see \cite{Manin}).

\vskip2mm
\noindent
Let $G$ be a connected algebraic group defined over $k$. The set of points in $G(k)$ that are $R$-equivalent to identity in $G(k)$ forms a normal subgroup of $G(k)$, denoted by $RG(k)$. The set $G(k)/R$ of $R$-equivalence classes in $G(k)$ is in canonical bijection with $G(k)/RG(k)$ and thus has a natural group structure. We identify $G(k)/R$ with the group $G(k)/RG(k)$. In the literature, it is known that this group is very useful in studying the rationality properties of $G$, for example, one may look at the references (\cite{Voskrenskii}, \cite{Gille1}).

\vskip2mm
\noindent
We call $G$ is \emph{$R$-trivial}, if $G(L)/R = \{1\}$ for all field extensions $L$ of $k$. It is a fact that if $G$ is $k$-rational, then $G$ is $R$-trivial (see
\cite{Voskrenskii}, Chap. $6$, Prop. $2$). Also, because tori of rank at most $2$ are rational, it follows that any reductive algebraic group of rank at most $2$ is rational (\cite{CM}, Cor. $2.3$). If $X$ and $Y$ are two irreducible varieties over $k$ then,
\begin{enumerate}
    \item $(X\times Y)(L)/R\cong X(L)/R\times Y(L)/R $.
    \vskip2mm
    \item $R_{F/k}(X)(k)/R=X(F)/R$.  (\cite{Voskrenskii})
\end{enumerate}

\vskip2mm
\noindent
If a group $G$ is $k$-split, then $G$ is $k$-rational (see \cite{VT}). It is also known that any connected unipotent algebraic group over a perfect field is rational (\cite{Voskrenskii}).

\vskip2mm
\begin{proposition}\label{Proposition 4.1}
Let $A$ be a finite-dimensional associative $k$-algebra. Assume that $\emph{\text{Aut}}_k(A)= \emph{Inn}(A)$. Then $\emph{\text{Aut}}_k(A)$ is $R$-trivial.
\end{proposition}

\vskip1mm
\begin{proof}
Consider the following rational map $\theta :\mathbb{A}^1_{k}\dashrightarrow \text{Inn}(A)$ defined by 
\begin{equation*}\label{eq4}
\theta(t)(x)=\{(1-t).1+ta\}x\{(1-t).1+ta\}^{-1}
\end{equation*}
\vskip1mm
\noindent
Therefore, $\theta(0)=Id$ and $\theta(1)=\text{Int}(a)$. It is clear that $L_a:A\rightarrow A$ is invertible if and only if $a$ is invertible, where $L_a(x)=ax$. 
%Suppose $L_a$ is invertible; then, as it is surjective, we have $ay=1$ for some $y\in A$. Consider $L_y$ as defined above, then $L_a\circ L_y=Id$, therefore, $L_y\circ L_a=Id$ this implies $ya=1$. For the converse part, suppose that $aa^{-1}=1$ then $L_{a^{-1}}$ serves as the inverse of $L_a$. 
It follows that $c=\{(1-t).1+t.a\}$ is invertible 
%if and only if $L_c$ is invertible. However, $L_c$ is invertible if and only if $\text{det}(L_c)\neq 0$, which implies that $c$ is invertible 
except for a finite set. Therefore, the map $\theta$ is well-defined. Hence, $k$-rational points are $R$-equivalent. Now, $\text{Aut}(A)(F)\cong \text{Aut}_F(A\otimes F)=\text{Inn}(A\otimes F)$ for any field extension $F$ over $k$. Again, by a similar line of argument, we can show that $F$-rational points are $R$-equivalent. Hence, we are done. 
\end{proof}


\vskip1mm

\begin{example}\label{Example 4.2}
Let $A=\bigoplus_{i=1}^{m} M_{n_i}(k)$, then by the classical result of Skolem and Noether $G_A\cong \prod_{i=1}^{m} \text{PGL}_{n_{i}}$. However, $\text{PGL}_n$ is $R$-trivial as a $k$-group. Hence, $G_A$ is $R$-trivial.
\end{example}


\vskip2mm

\begin{example}\label{Example 4.3}
Let $A$ be the set of upper triangular matrices in $M_n(k)$. Here, the radical $J=$ upper triangular nilpotent matrices and $A_s=$ diagonal matrices, and we have $A=A_s\oplus J$ as the Wedderburn-Malcev decomposition. Then $\text{Aut}_k(A)=\text{Inn}(A)$ as shown in (\cite{Pollack}, \cite{Stanley}). Therefore, this group is $R$-trivial by Proposition \ref{Proposition 4.1}.
\end{example}

\vskip2mm

\begin{example}\label{Example 4.4}
Let $V$ be an $n$-dimensional vector space over a field $k$ and let $A$ be the exterior algebra constructed from $V$. Then by Example $1.5$ in \cite{Pollack}, $\text{Aut}_k(A)\cong \text{GL}_n\ltimes U$, where $U$ is a unipotent group. Therefore, $\text{Aut}_k(A)$ is rational because unipotent groups are rational, hence $R$-trivial.
\end{example}

\vskip2mm

\begin{example}\label{Example 4.5}
Let $A=\displaystyle\frac{k[X_1,X_2,\dots,X_n]}{\langle X_1,X_2,\dots,X_n\rangle^l}$ for some $l\geq 2$, then $\text{Aut}_k(A)\cong \mathrm{GL}_n\ltimes U$, where $U$ is a unipotent group (see \cite{AS1}, Theorem $1.3)$. Therefore, $\text{Aut}_k(A)$ is rational, hence $R$-trivial.
\end{example}

\vskip2mm

\begin{example}\label{Example 4.6}
In \cite{AS1}, it has been proved that if $A$ is a split local commutative associative monomial algebra, then the identity component $G_A$ has rank $n$ with a split torus of rank $n$. So, as a $k$-group, $G_A$ is $R$-trivial as any $k$-split group is rational by \cite{VT}. 
\end{example}

\vskip2mm

\begin{proposition}\label{Proposition 4.7}
$($\cite{Voskrenskii}, \emph{Section} $17)$ If every maximal torus within a connected algebraic group \( G \) defined over a perfect field is \( R \)-trivial, then the group is \( R \)-trivial.
\end{proposition}

\vskip2mm

\begin{proposition}\label{Proposition 4.8}
$($\cite{CM}, \emph{Proposition} $2.2)$ If any maximal torus of a connected reductive group $G$ is rational (resp. stably rational) over its field of definition, then the group $G$ is itself rational (resp. stably rational). 
\end{proposition}
\vskip1mm
\begin{proof}
A similar proof will work for the stably rational case as done in Proposition $2.2$ of \cite{CM} for the rational case with the help of Section $4.1$ in \cite{Voskrenskii}.
\end{proof}
\vskip2mm
\begin{remark}\label{Remark 4.9}
For any connected algebraic group $G$ over a perfect field $k$, we have $G\cong_{b} G/U\times U$, where $U$ is the $k$-unipotent radical. Therefore, $G$ is rational (resp. stably rational) if and only if $G/U$ is rational (resp. stably rational), as $U$ is always rational. However, $G/U$ is a reductive group. Hence, if any maximal torus of $G$ is rational (resp. stably rational), then $G$ is also rational (resp. stably rational) by Proposition \ref{Proposition 4.8}.
\end{remark}
\vskip2mm
\begin{definition}\label{Definition 4.10}
A grading on an algebra $A$ is a family of vector subspaces $(A_n)_{n\in\mathbb{Z}}$ of $A$ such that $A=\bigoplus_{n\in\mathbb{Z}} A_n$ and $A_n.A_m\subseteq A_{n+m}$ for all $n,m\in \mathbb{Z}$. The grading is called \emph{by the radical} when $A_m=0$ for $m<0$ and $J^m=\bigoplus_{n\geq m} A_n$, for each $m\geq 0$ (see \cite{Saorin}, Proposition $1.1$ for equivalent definitions$)$.  For any associative algebra $A$ with $J^0=A$, we can associate a graded by the radical algebra, $G_{r}(A)\coloneq \bigoplus_{m\geq 0} J^m/J^{m+1}$ (\emph{associated graded by the radical algebra}), as a vector space, and where the multiplication $(a+J^{m+1}). (b+J^{n+1})=(ab+J^{m+n+1})$ for $a\in J^{m}$ and $b\in J^{n}$, extends to all of $G_{r}(A)$ by linearity. An algebra $A$ is said to be \emph{graded by the radical algebra} if it has a grading by the radical, or equivalently, $A$ is isomorphic to $G_{r} (A)$. For more details, we refer to \cite{AS1} and \cite{Saorin}.
\end{definition}
\vskip2mm
\begin{remark}\label{Remark 4.11}
It is easy to see that there are well-defined maps for a split local algebra $A$ over $k$, which are given below. Since for any automorphism $\phi\in \mathrm{Aut}_{k}(A)$, it induces an automorphism $\bar{\phi}:J/J^2\rightarrow J/J^2$. We have a morphism of algebraic groups $\Phi_A: \mathrm{Aut}_{k}(A)\rightarrow \mathrm{GL}(J/J^2)$, where $\Phi_A(\phi)=\bar{\phi}$. If $\phi\in \mathrm{Aut}_{k}(A)$, then we can define a $k$-linear bijective map $\Psi_A(\phi)$, which is just the direct sum of the maps induced by $\phi$ on all $J^m/J^{m+1}$, $m\in \mathbb{N}\cup \{0\}$. It is clear that $\Psi_A(\phi)$ is actually an algebra automorphism, thus yields a map $\Psi_A:\mathrm{Aut}_{k}(A)\rightarrow \mathrm{Aut}_{k}(G_{r}(A))$. Similarly, for every $m\in\{2,\dots,n\}$, where $J^{n+1}=0$, we have a canonical map $\phi_{k}: \mathrm{Aut}_{k}(A)\rightarrow \mathrm{Aut}_{k}(A/J^m)$ that takes $\phi\in \mathrm{Aut}_{k}(A)$ to the induced automorphism $\tilde{\phi}: A/J^m\rightarrow A/J^m$.  It is easy to see that we have similar kind of maps $\Phi_{G_{r}(A)}: \mathrm{Aut}_{k}(G_{r}(A))\rightarrow \mathrm{GL}(J/J^2)$ and $\Phi_{A/J^{m}}:\mathrm{Aut}_{k}(A/J^m)\rightarrow \mathrm{GL}(J/J^2)$.  All of the above maps $(\Phi_A, \Psi_A, \phi_{k},\Phi_{G_{r}(A)}, \Phi_{A/J^m}$) are morphism of algebraic groups.
\end{remark}
\vskip2mm
\noindent
The following result is known from \cite{Pollack}. We include a proof here for convenience.

\vskip2mm
\begin{proposition}\label{Proposition 4.12}
Let $A$ be a split local associatve algebra over $k$ and  $G_r(A)$ is defined as above, i.e., $G_{r}(A)\coloneq (k\oplus\bigoplus_{m=1}^{n}(J^m/J^{m+1}))$ be the associated graded by the radical algebra of $A$. Consider the following commutative diagrams $(1)$ and $(2)$:
\begin{equation}\label{eq6}
\begin{minipage}{0.45\linewidth}
\begin{tikzcd}
    \mathrm{Aut}_k(A) \arrow[dr, "\Phi_A"] \arrow[d, "\Psi_A"']\\
    \mathrm{Aut}_k(G_{r}(A)) \arrow[r,"\Phi_{G_{r}(A)}"'] & \mathrm{GL}(J/J^2)
\end{tikzcd}
\begin{center}
\textbf{(1)}
\end{center}
\end{minipage}
\begin{minipage}{0.45\linewidth}
\begin{tikzcd}
   \mathrm{Aut}_k(A) \arrow[dr, "\Phi_A"] \arrow[d, "\phi_k"']\\
    \mathrm{Aut}_k(A/J^m) \arrow[r,"\Phi_{A/J^m}"'] & \mathrm{GL}(J/J^2)
\end{tikzcd}
\begin{center}
\textbf{(2)}
\end{center}
\end{minipage}
\end{equation}
Then each of the induced morphisms $(\Phi_A, \Phi_{G_{r}(A)}, \Phi_{A/J^m})$ has a unipotent kernel. Therefore, the rank of $\emph{\text{Aut}}_k(A)$ is bounded by the dimension of $J/J^2$.
\end{proposition}
\vskip1mm
\begin{proof}
We will prove that $\text{Ker}(\Phi_A)$ is unipotent; other cases follow similarly. Since $A$ is a split local algebra, $\text{Aut}_k(A)=\text{Aut}_k(J)$. Therefore, \[\text{Ker}(\Phi_A)=\{\phi\in \text{Aut}_k(J):\phi(x)-x\in J^2, \forall x\in J\}=\{\phi\in \text{Aut}_k(J):(\phi -Id)(J)\subset J^2\}.\] Let $x\in J$ and $(\phi-Id)(x)=\sum_{i=1}^n x_iy_i$, where each $x_i,y_i\in J$. Then,
\begin{align*}
(\phi-Id)\bigl((\phi -Id)(x)\bigl)={} & 
\sum_{i=1}^n\phi(x_iy_i)-\sum_{i=1}^nx_iy_i
\end{align*}
which implies,

\begin{flalign*}
 (\phi-Id)^2(x) ={}&
  \!\begin{aligned}[t] 
    &\sum_{i=1}^{n}\phi(x_i)\phi(y_i)-x_1\phi(y_1)+x_1\phi(y_1)-\dots\\ 
     &-x_n\phi(y_n)+x_n\phi(y_n)-\sum_{i=1}^nx_iy_i
  \end{aligned}\\
  ={}&
  \!\begin{aligned}[t]
    &\sum_{i=1}^n(\phi-Id)(x_i)\phi(y_i)+\sum_{i=1}^nx_i(\phi-Id)(y_i) \\
  \end{aligned}\\
\end{flalign*}
\noindent
However, $(\phi-Id)(x)\in J^2, \forall x\in J$. Hence, $(\phi-Id)^2(J)\subset J^3.$ Suppose $J^m=0$ for some $m\geq 2$. So, we can iterate the process and show that $(\phi-Id)^{m-1}=0.$ Since the kernel is a unipotent group, any maximal torus of $\text{Aut}_k(A)$ embeds in $\text{GL}(J/J^2)$. Hence, the rank of any maximal torus is less than or equal to $\text{dim}(J/J^2)$.
 \end{proof}
 \vskip2mm
 \begin{remark}\label{Remark 4.13}
Restricting $\Phi_A$ to $ G_A$ yields a homomorphism from $ G_A$ to $\text{GL} (J/J^2)$. By a similar line of argument, we can prove that $\text{Ker}(\Phi_A|_{G_{A}})$ is a unipotent group. Similarly, we can apply restrictions to other mappings to their identity components. Still, the proposition remains valid. To avoid notational ambiguity, the notation used for the identity component mappings will be the same as $\Phi_A$ and $\Psi_A$.
 \end{remark}

\vskip2mm

\begin{proposition}\label{Proposition 4.14}
Let $A$ be a split local associative algebra over $k$ such that $\emph{dim}(J/J^2)=\emph{rank}(G_A)=n$, then $G_A$ is rational.
\end{proposition}
\vskip1mm
\begin{proof}
We have the canonical map 
\begin{equation*}
\Phi_A: G_A \rightarrow \text{GL}(J/J^2)
\end{equation*}
whose kernel is unipotent by Proposition \ref{Proposition 4.12}. Therefore, any maximal torus of $G_A$ is isomorphic to a maximal torus of $\text{GL}(J/J^2)$. However, we know that any maximal torus of $\text{GL}_n$ over $k$ is isomorphic to \[R_{L_{1}/k}(\mathbb{G}_{m, L_1})\times R_{L_{2/k}}(\mathbb{G}_{m, L_{2}})\times\dots \times R_{L_{s}/k}(\mathbb{G}_{m,L_s}),\] where $L_i$'s are separable extensions of $k$ and $\sum_{i=1}^s[L_i:k]=n$. It follows from Section $3.12$ of \cite{Voskrenskii} that the Weil restriction of $\mathbb{G}_m$ is rational. Hence, the result follows from the Proposition \ref{Proposition 4.8} and Remark \ref{Remark 4.9}.
\end{proof}
\vskip2mm

\begin{proposition}\label{Proposition 4.15}
Let $A$ be a split local associative algebra satisfying $\emph{dim}(J/J^2)\leq 5$, then $G_A$ is stably rational, hence $R$-trivial.
\end{proposition}
\vskip1mm
\begin{proof}
Let $\Phi_A$ be the canonical map $\Phi_A:G_A\rightarrow \text{GL}(J/J^2)$ and consider 
\begin{equation*}\label{eq8}
\text{GL}(J/J^2)\xrightarrow{\text{det}} \mathbb{G}_m 
\end{equation*} 
If the $\text{dim}(J/J^2)$ is $1$ or $2$, then the result follows from the fact that any torus of rank at most $2$ is rational.
\vskip2mm
\noindent
\textbf{Case $1$, $\text{dim}(J/J^2)=3$:} Let $T$ be any maximal torus of $G_A$ that embeds in $\text{GL}(J/J^2)$. Therefore, the rank of $T$ is less than or equal to $3$. If $\text{det}(T)=1$, then $T\subset \text{SL}(J/J^2)$, and thus dimension of $T$ is less than or equal to $2$. On the other hand, if $\text{det}(T)\neq 1$, then $T$ has a non-trivial character namely, $\text{det}|_{T}$. Therefore, $T\cong_{b} \mathbb{G}_m^{r}\times_k T_a$, where $r\geq 1$ and the dimension of the anisotropic part $T_a$ of the torus $T$ is less than or equal to $2$ by Section $17.3$ of \cite{Voskrenskii}. Hence, in this case, using rationality of rank $\leq 2$ tori, it follows that $T$ is rational. 
\vskip2mm 
\noindent
\textbf{Case $2$, $\text{dim}(J/J^2)=4$:} In this case the possible rank of a maximal torus $T$ in $G_A$ is $1,2,3,4$. If the rank of $T$ is $ 1,2\text { or } 4$, then it follows from the rationality of tori rank $\leq 2$ and Proposition \ref{Proposition 4.14} that $T$ is rational. Suppose $\text{rank}(T)=3$ and $\text{det}(T)$ is non-trivial, then we are done by an argument similar to the previous case. Suppose $\text{det}(T)=1$, then $T\subset \text{SL}(J/J^2)$. Therefore, $S\coloneq \mathbb{G}_{m}.T$ is a maximal torus in $\text{GL}_{4}$, where $\mathbb{G}_m$ is the torus of scalar matrices. If $T$ has a non-trivial character defined over $k$, then $T\cong_b T_a\times \mathbb{G}_m^{r}$, $r\geq 1$ and $T_a$ is anisotropic of rank $\leq 2$. From this, it follows that $T$ is rational. Hence we may assume that $S\cong_b S_a \times \mathbb{G}_m^{r}$($r\geq 1$) with $S_a=T$, where $S_a$ is the anisotropic part of $S$. Again, we know that $S$ is rational; therefore, $T$ is a stably rational torus. Hence, we have shown that any maximal torus of $G_A$ under the given hypothesis is either rational or stably rational. As any rational torus is stably rational, therefore, Proposition \ref{Proposition 4.8} and Remark \ref{Remark 4.9} complete the proof.
\vskip2mm
\noindent
\textbf{Case $3$, $\text{dim}(J/J^2)=5$:} The rank of the maximal torus in this case is possibly $1,2,3,4,5$. If it is $1,2,5$ then we are done by the rationality of tori rank $\leq 2$ and Proposition \ref{Proposition 4.14}. Suppose, $\text{rank}(T)=3\text{ or }4$ and $\text{det}(T)\neq 1$, then $T$ is isotropic and $T\cong_b \mathbb{G}_m^{r}\times T_a$, $r\geq 1$ and the rank($T_a$) is at most $2 \text{ or } 3$, and $T_a\subset \text{SL}(J/J^2)$. If $\text{rank}(T_a)\leq 2$, $T_a$ is rational. If $\text{ det}(T)=1$ and $\text{rank}(T)=4$, then $T$ is contained in $\text{SL}_5$. It follows that $S\coloneq \mathbb{G}_m.T$ is a maximal torus in $\text{GL}_5$, and since $S$ is rational and $S\cong_b S_a\times \mathbb{G}_m^r$, $r\geq 1$; as before, $T$ is stably rational. 
\vskip2mm
\noindent
So, in both the remaining cases, it reduces to the case where $T$ is anisotropic and $\text{rank}(T)=3$.  Assume now $T$ is anisotropic and $\text{rank}(T)=3$. Since $T\subset \text{SL}(J/J^2)$, $T\subset L^{(1)}$, the norm $1$ torus of an \'etale algebra over $k$ of degree $5$. Moreover, there exists a rank $1$ torus $T'$ such that $T.T'=L^{(1)}$ by Proposition $13.2.3$ of \cite{Springer}. If $T'=\mathbb{G}_m$, then we can write $L^{(1)}\cong_b T\times \mathbb{G}_m$. However, we know that $S=L^{(1)}.\mathbb{G}_m$ is a maximal torus in $\text{GL}_5$, where $\mathbb{G}_m$ is the torus of scalar matices. It implies that $S\cong_b (T\times\mathbb{G}_m).\mathbb{G}_m$. Since $T$ is anisotropic, we have $S\cong_b S_{a}\times \mathbb{G}_m^2$ with $S_a=T$. Therefore, $T$ is stably rational. Suppose $T'=K^{(1)}$, where $K$ is a degree $2$ extension of $k$, then 
\begin{equation}\label{eq 4.0.7}
T.K^{(1)}=L^{(1)}.
\end{equation}
Since the left-hand side of the previous equation \ref{eq 4.0.7} is anisotropic, $L^{(1)}$ must be anisotropic. Hence, $L$ has only two choices: either $L$ has to be a field extension of degree $5$ over $k$ or $L=E\times F$, where $E$ is a degree $3$ field extension and $F$ is a degree $2$ field extension over $k$. For other possibilities of \'etale algebras $L$ give a norm $1$ torus which is isotropic, which implies $\mathbb{G}_{m,k}$ embeds in $T$ as $\mathbb{G}_{m,k}\cap K^{(1)}=\{1\}$, which is a contradiction as $T$ is anisotropic. For example if $L=k\times F_1\times F_2$, where $F_1$ and $F_2$ both are field extension of degree $2$ over $k$, then \[L^{(1)}=\{(a,b,c)\in k\times F_1\times F_2:N(a)N(b)N(c)=1\}, \text{ where $N$ is the norm map, }\] which implies $L^{(1)}=\{(1/N(b)N(c),b,c):b\in F_{1}^{\times}, c\in F_{2}^{\times}\}$. Hence $L^{(1)}$ is isotropic. If $L$ is a field extension of degree $5$ over $k$, then $L^{(1)}$ is isotropic over $K$ as $K^{(1)}$ is isotropic over $K$, and it contains $K^{(1)}$ as $L^{(1)}=T.K^{(1)}$, which is not possible because $[L:k]=5$ and $[K:k]=2$. So we have $T.K^{(1)}=(E\times F)^{(1)}$. However, $T.K^{(1)}$ is isotropic over $K$. Therefore, $(E\times F)^{(1)}$ is isotropic over $K$ and contains $K^{(1)}$. Hence, $K\cap F\neq k$ and $[K:k]=[F:k]=2$ imply $K\cap F=F=K$. So finally we get that $T.K^{(1)}=(E\times K)^{(1)}$. It is easy to see that $E^{(1)}\times K^{(1)}\subset TK^{(1)}$. If $e\in E^{(1)}$, then $e=t\alpha$ with $t\in T$ and $\alpha\in K^{(1)}$. Suppose $\alpha =1$ for all $e\in E^{(1)}$, then we have $E^{(1)}\subset T$ which implies $T\cong E^{(1)}.K'^{(1)}$, where $K'$ is a field extension of degree $2$ over $k$ as $T$ is anisotropic. However, $E^{(1)}\cap K'^{(1)}=\{1\}$ as $[E:k]=3$ and $[K':k]=2$. Therefore $T\cong E^{(1)}\times K'^{(1)}$. Hence, $T$ is rational as both $E^{(1)}$ and $K'^{(1)}$ are rational. If $\alpha\neq 1$ for some $e\in E^{(1)}$, then $t=e\alpha^{-1}$ and the Zariski closure of the cyclic group generated by ${\{(e,\alpha^{-1})\}}$ is the group $E^{(1)}\times K^{(1)}\subset T$. It implies that $T=E^{(1)}\times K^{(1)}$. Therefore, $T$ is rational. Hence, it concludes the proof for $\text{dim}(J/J^2)=5$. In all the above cases, we have proved that any maximal torus in $G_A$ is stably rational. Therefore, the result follows from Proposition \ref{Proposition 4.8} and Remark \ref{Remark 4.9}.
\end{proof}


\vskip2mm

\begin{corollary}\label{Corollary 4.16}
Let $A$ be a split local algebra with $\mathrm{dim}_k A\leq 7$. Then, $G_A$ is stably rational, hence $R$-trivial.
\end{corollary}
\vskip1mm
 \begin{proof}
 If $J^2=0$, then $\text{Aut}_k(A)=\text{Aut}_k(J)=\text{GL}(J)$. So, $G_A$ is rational, hence $R$-trivial. Suppose $J^2\neq0$. Since $A$ is a local algebra with $\text{dim }A\leq 7$, $\text{dim }J\leq 6$ and $\text{dim }J^2\leq 5$ (if $J^2=J$ then $J=0$). Hence, we get $\text{dim}(J/J^2)\leq 5$. Using Proposition \ref{Proposition 4.15}, we see that $G_A$ is stably rational.
\end{proof}

\vskip2mm

\begin{remark}\label{Remark 4.17}
For any split local algebra $A$, if a maximal torus in $\text{Aut}_k(A)$ has dimension equal to the dimension of $J/J^2$, it must be isotropic. Let $T$ be a maximal torus in $G_A$, then $\text{det}(T)\neq 1$, otherwise $T\subset \text{SL}(J/J^2)$ which implies that $\text{dim}(T)<\text{dim}(J/J^2)$. Therefore $T$ has a non-trivial character defined over $k$, namely $\text{det}|_{T}$.
\end{remark}

\vspace{0.2in}
\noindent
We will often need the following result:
\vskip2mm
\begin{lemma}\label{Lemma 4.18}
$($\cite{CP}, \emph{Lemma} $1)$ Let $G$ be a connected algebraic group defined over an arbitrary field $k$. If $H$ is a connected subgroup of $G$ defined over $k$ such that for any field extension $F/k$, the kernel of the natural morphism $H^1(F,H)\rightarrow H^1(F,G)$ is trivial, then the variety 
\begin{equation*}\label{eq9}
G\cong_{b} G/H \times H.\end{equation*}
\end{lemma}
\vskip2mm
\begin{remark}\label{Remark 4.19}
Let $G$ be a connected algebraic group over a perfect field $k$ and $U$ be a $k$-unipotent subgroup of $G$, then $G\cong_b G/U\times U$ by using Lemma \ref{Lemma 4.18} as $H^{1}(F, U)$ is trivial for the unipotent group $U$ for every field extension $F$ over $k$.  
\end{remark}
\vspace{0.2in}

\noindent
The above results motivate us to ask whether every connected algebraic group of the form $G_A$, where $A$ is a finite-dimensional associative algebra over a field $k$, is $R$-trivial or not. Below, we prove a key result on $R$-triviality of $G_A$.

\vspace{0.2in}
\begin{theorem}\label{Theorem 4.20}
Let $A$ be a split finite-dimensional associative algebra over a perfect field $k$ with $A_s$ a semisimple subalgebra of $A$ such that $A=A_s\oplus J$ (vector space decomposition) and $G_{A, A_s}\coloneq\{\sigma\in G_A: \sigma(a)=a \quad \forall a\in A_s\}$. Then, $G_A$ is $R$-trivial if and only if $G_{A,A_s}$ is $R$-trivial. Suppose any one of the following holds: 
\begin{enumerate}
    \item $J^2=0$;
    \vskip2mm
    \item $\emph{dim}(J/J^2)\leq 5$;
    \vskip2mm
    \item $\emph{rank}(G_{A,A_s})=\emph{dim}(J/J^2)$;
\end{enumerate}
\vskip2mm
 Then $\emph{\text{Aut}}_k(A)^{0}\coloneqq G_A$ is $R$-trivial.
\end{theorem}
\vskip1mm
\begin{proof}
Using the decomposition from Corollary \ref{Corollary 2.2}, we can write $G_A$ as a product of some useful subgroups as Pollack did in \cite{Pollack}. Under the given hypothesis, we prove that these groups are all $R$-trivial. First, we prove some lemmas to decompose the group $G_A$.

\vskip2mm
\begin{lemma}\label{Lemma 4.21}$($\cite{Pollack}$)$
$G_A=\hat{J}G_{A}^{A_s}$, where $\hat{J}=\{\emph{Int}(1+r)\in G_A : r\in J\},\text{ and } G_{A}^{A_s}\coloneq\{\sigma \in G_A: \sigma(A_s)=A_s\}.$
\end{lemma}

\vskip1mm
\begin{proof}
Let $\sigma \in G_A$, then $A_s\oplus J=A=\sigma(A)=\sigma(A_s)\oplus J$. Therefore, using Corollary \ref{Corollary 2.2} we have $\sigma(A_s)=\rho(A_s)$, for some $\rho \in \hat{J}$. Thus, $\rho^{-1}\sigma\in G_{A}^{A_{s}}$ and it follows that $G_A=\hat{J}G_{A}^{A_s}$. Hence, \begin{equation*}\label{eq10}
\displaystyle\frac{G_A}{\hat{J}}\cong \displaystyle\frac{G_{A}^{A_s}}{G_{A}^{A_s}\cap \hat{J}}
\end{equation*}
\end{proof}

\vskip2mm

\begin{remark}\label{Remark 4.22}$($\cite{Pollack}$)$
If $J^2=0$, then for $G_A=\hat{J}G_{A}^{A_s}$, the product is a semidirect product.
\end{remark}

\vskip4mm
\noindent
Note that $\text{Int}(1+J)=\hat{J}$ is unipotent (rational). It is clear from Lemma \ref{Lemma 4.18} and Remark \ref{Remark 4.19} that $G_A$ is $R$-trivial if and only if $G_A/\hat{J}$ is $R$-trivial. Hence, $G_A$ is $R$-trivial if and only if $\displaystyle\frac{G_{A}^{A_s}}{G_{A}^{A_s}\cap \hat{J}}$ is $R$-trivial. Again using Lemma \ref{Lemma 4.18}, we can conclude that $G_A$ is $R$-trivial if and only if $G_{A}^{A_s}$ is $R$-trivial. Therefore, it is enough to prove that $G_{A}^{A_s}$ is $R$-trivial.

\vskip4mm
\begin{lemma}\label{Lemma 4.23}
$G_{A}^{A_s}\cong \emph{Inn}(A_{s}^{\times})G_{A,A_s}$, where $\emph{Inn}(A_s^{\times})\coloneq\{\emph{Int}(x)\in G_{A/J} : x\in A_{s}^{\times}\}$ and $G_{A,A_s}=\{\sigma\in G_A: \sigma|_{A_s}=Id\}$ and $A_s^{\times}=\emph{Units of }A_{s}$.
\end{lemma}
\vskip1mm
\begin{proof}
Consider the induced map 
\begin{equation*}\label{eq11}
G_A\rightarrow \text{Aut}_k(A/J)
\end{equation*}
whose image is closed and connected and is contained in $\text{Inn}(A/J)=G_{A/J}$ as $A/J\cong \bigoplus_i M_{n_i}(k)$ . Therefore, we have an exact sequence 
\begin{center}
\begin{tikzcd}
1\arrow[r] & G_{A,A_s} \arrow[r] &G_{A}^{A_s} \arrow[r,"res"] & \text{Inn}(A_s^{\times})\arrow[bend left=33]{l}{s}\arrow[r] & 1,
\end{tikzcd}
\end{center}
\vskip1mm
\noindent
which has a section $s:\text{Inn}(A_s^{\times})\rightarrow G_{A}^{A_s}$ such that $s(\text{Int}(a)_{|A_s})=\text{Int}(a)$, where $a\in A_s^{\times}$, $res(\phi)=\phi_{|A_s}$, $\phi\in G_{A}^{A_s}$. Therefore, the above sequence is split exact. Hence, the result follows.
\end{proof}

\vskip2mm  
\noindent
So, by the previous lemma to prove $G_{A}^{A_s}$ is $R$-trivial, it suffices to prove that $G_{A, A_s}$ is $R$-trivial, as in Proposition \ref{Proposition 4.1} we have already shown that the automorphism groups whose elements are inner automorphisms are $R$-trivial. Hence, it follows that $G_A$ is $R$-trivial if and only in $G_{A, A_S}$ is $R$-trivial.
\vskip4mm
\noindent
Since $J$ is an ideal of $A$, it may be considered as a $A_s$-bimodule or equivalently, as a left module for the semisimple algebra $A_{s}^{\#}\coloneqq A_s\otimes A_{s}^{\text{}op}$, where $A_{s}^{\text{}op}$ is the opposite algebra of $A_s$. Therefore, $J$ is an $A_{s}^{\#}$-algebra. We have for any $\sigma\in G_{A,A_s}$, $\sigma(arb)=a\sigma(r)b$, where $a,b\in A_s$ and $r\in J$. So, $G_{A,A_s}$ acts as an $A_s^{\#}$-algebra automorphism of $J$, and we have an injection $\Gamma: G_{A,A_s}\longrightarrow \text{Aut}_{A_s^{\#}-\textit{algebra}}(J)$. 

\vskip4mm
\begin{lemma}\label{Lemma 4.24}
$G_{A,A_s}\cong \emph{\text{Aut}}_{A_s^{\#}-\textit{algebra}}(J)$.
\end{lemma}
 \vskip1mm
\begin{proof}
It is clear from the previous argument that we have an injection 
\begin{equation*}\label{eq13}
\Gamma: G_{A,A_s}\longrightarrow \text{Aut}_{A_s^{\#}-\textit{algebra}}(J)
\end{equation*}
However, for $\psi\in \text{Aut}_{A_s^{\#}-\textit{algebra}}(J)$,  define $\bar \psi(a+r)=a+\psi(r)$, where $a\in A_s$ and $r\in J$, then $\Gamma(\bar{\psi})=\psi$ and the above map is surjective. Since $\Gamma(\psi_1\circ\psi_2)=\Gamma(\psi_1)\Gamma(\psi_2)$, we have the isomorphism.
\end{proof}

\vskip2mm
\noindent
So, it is enough to prove that $\text{Aut}_{A_{s}^{\#}-\text{algebra}}(J)$ is $R$-trivial. Hence, the $R$-triviality of $G_A$ depends on the radical of $A$. Since $A$ is an artinian algebra, there must exist an integer $n$ such that $J^n=0$. Next, we consider the case $n=2$.

\vskip2mm
\begin{lemma}\label{Lemma 4.25}
If $J^2=0$, then $G_{A,A_s}\cong \emph{\text{Aut}}_{A_s^{\#}-\textit{module}}(J)$ and it is a rational group.
\end{lemma}
\vskip1mm
\begin{proof}
For $J^2=0$, we have $\text{Aut}_{A_{s}^{\#}-\textit{algebra}}(J)\cong \text{Aut}_{A_s^{\#}-\textit{module}}(J)$, where $J$ is a $A_s^{\#}$-module. It follows from the discussion of Chapter $2$ of \cite{Tifr} and \cite{Pierce} that the $A_s^{\#}$-\text{module} $J$ can be written as 
$J=\sum_{i=1}^{\alpha} \lambda_{i}M_{i}$,
where this $M_i$'s are inequivalent irreducible $A_s^{\#}$-modules and $\lambda_i\in \mathbb{N}$. By using Schur's lemma it follows that $\text{Aut}_{A_s^{\#}-\text{module}}(\lambda_i M_i)=\text{GL}_{\lambda_i}$. But
\begin{equation*}\label{eq15}
\text{Aut}_{A_s^{\#}-\textit{module}}(J)\cong \prod_{i=1}^{\alpha}\text{Aut}(\lambda_{i}M_i), \text{ where $\prod$ denotes the product of groups. }
\end{equation*}
Therefore, we get \[G_{A,A_s}\cong \prod_{i=1}^{\alpha} \text{GL}_{\lambda_i}.\]

\vskip2mm
\noindent
Since $\text{GL}_n$ is rational, and a product of rational groups is rational, it follows that $G_{A, A_s}$ is rational. Therefore, $k$-rational points of $\text{Aut}_k(A)^{0}$ are $R$-equivalent over $k$. If we consider the $F$-rational points $\text{Aut}_k(A)^{0}(F)=\text{Aut}_F(A\otimes F)^{0}$, where $F$ is an extension of $k$, then the assertion follows easily by using Theorem \ref{Theorem 3.5} and Remark \ref{Remark 3.7}. Hence, $F$-rational points are also $R$-equivalent. Therefore, $\text{Aut}_k(A)^{0}=G_A$ is $R$-trivial.
\end{proof}

\vskip2mm

\begin{lemma}\label{Lemma 4.26}
If $J^2\neq 0$ and $\emph{dim}(J/J^2)\leq 5$ or $\emph{rank}(G_{A,A_s})=\emph{dim}(J/J^2)$, then $G_{A,A_s}$ is $R$-trivial.
\end{lemma}
\vskip1mm
\begin{proof} Since $J^2\neq 0$ and $G_{A,A_s}\cong \text{Aut}_{{A_s^{\#}}-\text{algebra}}(J)$, we have the canonical homomorphism \[G_{A,A_s}\rightarrow \text{Aut}_k(J/J^2).\] Now using similar argument as in Proposition \ref{Proposition 4.12}, we can show that any maximal torus of $G_{A,A_s}$ embeds in $\text{Aut}_k(J/J^2)=\text{GL}(J/J^2)$. Therefore, by using a similar line of argument as done in Proposition \ref{Proposition 4.15}, we can show that any maximal torus of $G_{A, A_s}$ is stably rational if $\text{dim}(J/J^2)\leq 5$. Hence, $G_{A, A_s}$ is stably rational if $\text{dim}(J/J^2)\leq 5$. If $\text{rank}(G_{A,A_s})=\text{dim}(J/J^2)$, then by Proposition \ref{Proposition 4.14}, $G_{A,A_s}$ is rational. Hence in any case $G_{A,A_s}$ is $R$-trivial.
\end{proof}
\vskip2mm
\noindent
Finally, we use Lemma \ref{Lemma 4.26} and Lemma \ref{Lemma 4.25}  to conclude that $G_A$ is $R$-trivial when it satisfies the conditions given in Theorem \ref{Theorem 4.20}.
\end{proof}

\vskip2mm
\begin{corollary}\label{Corollary 4.27}
Let $\emph{\text{Aut}}_k(A)$ be a connected reductive group, where $A$ is a split finite-dimensional associative algebra, then $\mathrm{Aut}_k(A)$ is $R$-trivial.
\end{corollary}
\vskip1mm
\begin{proof}
For any connected reductive group of the form $\text{Aut}_k(A)$, we have $J^2=0$ by Corollary \ref{Corollary 3.2}, where $J$ is the radical of $A$. It follows from Theorem \ref{Theorem 4.20} that $\text{Aut}_k(A)$ is $R$-trivial.
\end{proof}

\vskip2mm
\begin{corollary}\label{Corollary 4.28}
Let $T$ be a non $R$-trivial torus over $k$. Then, there does not exist any split finite-dimensional associative algebra $A$ over $k$ such that $\emph{\text{Aut}}_k(A)\cong T$.
\end{corollary}
\vskip 1mm
\begin{proof}
Follows from Corollary \ref{Corollary 4.27}.
\end{proof}
\vskip2mm
\begin{corollary}\label{Corollary 4.29}
Let $A$ be a split finite-dimensional associative algebra over $k$. For every field extension $F$ over $k$, let us define $R_{G_A}(F)$, \[R_{G_A}(F)\coloneq\{\sigma\in G_A(F) : \sigma \text{ is R-equivalent to identity } \} .\] Then, $R_{G_A}(F)=\{Id\}$ for all field extension $F$ over $k$ if and only if $A\cong \bigoplus_i k$.
\end{corollary}
\vskip1mm
\begin{proof}
It is already known that $R_{G_A}(F)$ is a normal subgroup of $G_A(F)$. Moreover, it is clear that any unipotent element is $R$-equivalent to the identity. Therefore, for every field extension $F$ over $k$, we have $U(F)\subset R_{G_A}(F)$, where $U$ is the $k$-unipotent radical of $G_A$. Suppose $R_G(F)=\{Id\}$ for all field extension $F$ over $k$, then $k$-unipotent radical of $U$ of $G_A$ is trivial. Hence, $G_A$ is reductive over the perfect field $k$. However, by Corollary \ref{Corollary 4.27}, $G_A$ is $R$-trivial, which means $G_A(F)=R_{G_A}(F)$ for every field extension $F$ over $k$. Hence, we get that $G_A=\{Id\}$ which further implies that $A\cong \bigoplus_i k$. This completes one direction of the proof, and the other direction follows easily.
\end{proof}
\vskip2mm
\noindent
It is easy to observe that using the proof of Theorem \ref{Theorem 4.20}, we can classify the connected reductive groups which are of the form $\text{Aut}_k(A)$, where $A$ is a split finite-dimensional associative algebra.

\vskip2mm

\begin{corollary}\label{Corollary 4.30}
Let $A$ be a split finite-dimensional associative algebra over the perfect field $k$ such that $\emph{\text{Aut}}_k(A)$ is a connected reductive group. Then \[\emph{\text{Aut}}_k(A)\cong (\emph{Inn}(A_s^{\times}).(\prod_{i=1}^{\alpha} \emph{GL}_{n_{i}})),\] for some $\alpha$, where $A_s$ is a semisimple subalgebra $A$ such that $A= A_s\oplus J$.
\end{corollary}
\vskip1mm
\begin{proof}
 Since $\text{Aut}_k(A)$ is a connected reductive group, we have $J^2=0$, where $J$ is the radical of $A$. Now using Remark \ref{Remark 4.22} we have $G_A=G_{A}^{A_s}\ltimes \hat{J}$, where $\hat{J}=\{\text{Int}(1+r):r\in J\}$. It follows from Lemma \ref{Lemma 4.23} that $G_{A}^{A_s}\cong\text{Inn}(A_s^{\times})G_{A,A_s}$. Since $\text{Aut}_{k}(A)=G_A$ is reductive, we have $\hat{J}=\{Id\}$. However, the proof of Lemma \ref{Lemma 4.25} shows that $G_{A,A_s}\cong \bigoplus_{i=1}^{\alpha} \text{GL}_{n_i}$. Hence, we get that
\begin{equation*}\label{eq17}
\text{Aut}_k(A)\cong (\text{Inn}(A_s^{\times}).(\prod_{i=1}^{\alpha} \text{GL}_{n_{i}}))
\end{equation*}
\end{proof}

\vskip2mm 
\noindent
The following corollary presents a remarkable insight.

\vskip2mm

\begin{corollary}\label{Corollary 4.31}
Let $A$ be a split finite-dimensional associative commutative algebra. If $\mathrm{Aut}_k(A)$ is a connected reductive group of rank $n$, then
\begin{equation*}\label{eq18}
\emph{\text{Aut}}_k(A)\cong \prod_{i}^{k} \emph{GL}_{n_{i}}
\end{equation*}
\text{ where } $\sum_{i}^{k} n_i=n.$
\end{corollary}
\vskip1mm
\begin{proof}
$\text {Inn}(A_s^{\times})=\{Id\}$ as $A$ is commutative. Therefore, using Corollary \ref{Corollary 4.30}, we are done.
\end{proof}

\vskip2mm
\noindent
To state our next result, we need the following lemma.

\vskip2mm

\begin{lemma}\label{Lemma 4.32}
$($\cite{Pollack}, \emph{Theorem} $1.6)$ Let $A$ be a split local associative algebra over a perfect filed $k$ such that $G_A$ is nilpotent. Then, either $G_A$ is a torus or unipotent. It can not be a direct product of a torus with a non-trivial unipotent group.   
\end{lemma}

\vskip2mm

\begin{theorem}\label{Theorem 4.33}
Let $A$ be a split finite-dimensional associative algebra over a perfect field $k$. Assume that $\emph{\text{Aut}}_k(A)$ is a nilpotent group. Then, $\emph{\text{Aut}}_k(A)^{0}\coloneqq  G_A$ is rational.
\end{theorem}
\vskip1mm
\begin{proof}
Since $\text{Aut}_k(A)$ is nilpotent, it is a solvable group, which further implies that $\text{Inn}(A_s^{\times})$ is solvable, where $A_s$ is a semisimple subalgebra $A$ with $A=A_s\oplus J$. Hence, we get that $A_s^{\times}$ is solvable. However, $A_s\cong \bigoplus_{i=1}^{m}M_{n_i}(k)$ and $A_s^{\times}$ solvable implies $A_s\cong \bigoplus_{i}^{m} k$. So, we can decompose $A\cong \bigoplus_{i}^{m} A_i$, where each $A_i$ is a split local algebra, and this decomposition is unique (see \cite{Auslander}, \cite{Pierce}).
It follows that $G_A\cong\prod_{i}^{m}G_{A_{i}}$. Hence, to prove the assertion, it is enough to prove that $G_A$ is rational when $A$ is a split local algebra and $G_{A}$ is a nilpotent group. A connected nilpotent algebraic group is a direct product of a torus with an unipotent group. However, using the previous Lemma \ref{Lemma 4.32}, we have either $G_A=T$, where $T$ is the torus, or $G_A$ is unipotent. If $G_A$ is a torus, then by Corollary \ref{Corollary 4.28} and Theorem \ref{Theorem 3.4}, $G_A$ is $R$-trivial as well as rational. On the other hand, if $G_A$ is an unipotent group, then it is rational as a variety. Therefore, the result follows.
\end{proof}

\vskip2mm
\begin{proposition}\label{Proposition 4.34}
For a split finite-dimensional associative algebra $A$, the $R$-triviality of $G_{A}$ is Morita invariant.
\end{proposition}
\vskip1mm
\begin{proof}
Corollary $19$ in \cite{AS3} shows that $G_{A,A_s}/(G_{A,A_s}\cap \hat{J})$ is a Morita invariant, where $\hat{J}=\{\text{Int}(1+r):r\in J\}$. In Theorem \ref{Theorem 4.20}, we have shown that $G_A$ is $R$-trivial if and only if $G_{A, A_{s}}$ is $R$-trivial. Since $G_{A,A_s}\cap \hat{J}$ is a unipotent group, it follows from Lemma \ref{Lemma 4.18} that, \[G_{A,A_s}\cong_{b} \{G_{A,A_s}/(G_{A,A_s}\cap \hat{J})\} \times (G_{A,A_s}\cap \hat{J}).\] Therefore, it is clear that $G_A$ is $R$-trivial if and only if $G_{A,A_s}/(G_{A,A_s}\cap \hat{J})$ is $R$-trivial. Hence, the $R$-triviality of $G_{A}$ is Morita invariant. That means if $A$ and $B$ are two Morita equivalent algebras over $k$, then $G_{A}$ is $R$-trivial if and only if $G_{B}$ is $R$-trivial. 
 \end{proof}

\vskip2mm
\begin{remark}\label{Remark 4.35}
Theorem \ref{Theorem 4.20} shows that the $R$-triviality of $G_A$ depends on $G_{A, A_s}$ for a split finite-dimensional associative algebra. However, from Proposition \ref{Proposition 4.34}, it is clear that if $A$ and $B$ are two unital split finite-dimensional associative algebras which are Morita equivalent, then the $R$-triviality of $G_{A}$ implies the $R$-triviality of $G_{B}$ and vice versa. It follows from the Wedderburn-Malcev decomposition that any split finite-dimensional unital associative algebra $A$ is isomorphic to $A_s \oplus J$, where $A_s \cong \bigoplus_{i=1}^{m}M_{n_{i}}(k)$ and $J$ is the radical of $A$. Again, from Theorem \ref{Theorem 2.5} and Section $6.6$ of \cite{Pierce}, we know that for any such algebra $A$, there exists a Morita equivalent algebra $B=B_s\oplus J^{\prime}$, where $B_s\cong \bigoplus_{i=1}^{m} k$ and $J^{\prime}$ is the radical of $B$. So $B\cong \bigoplus_{i=1}^{m} B_i$, where each $B_i$ is a split local associative algebra. Now, if we consider the identity component of the automorphism group of $B$, we get $G_B\cong\prod_{i=1}^{m} G_{B_{i}}$. So, to study the $R$-equivalence property of the automorphism group of a split finite-dimensional associative algebra, one possible approach is to discuss the $R$-equivalence property of the automorphism group of split local associative algebras.
\end{remark}
\vskip2mm 
\noindent 
Next, we recall a result from \cite{Merkurjev}.
\vskip2mm

\begin{lemma}\label{Lemma 4.36}$($\cite{Merkurjev}, \emph{Lemma} $2.1)$
Let $f: Y\rightarrow X$ be a surjective morphism of varieties over $k$. Suppose that for every field extension $K/k$, and every point $x\in X(K)$, the fiber of $f$ over $x$ is a rational variety. Then $Y$ is stably birationally equivalent to $X$.
\end{lemma}
 
\vskip2mm
\begin{proposition}\label{Proposition 4.37}
Let $A$ be a split local algebra, then $G_A$ is stably birationally equivalent to $\emph{\text{Im}}(\Phi_A)$, where $\Phi_A$ is the canonical map, defined in Proposition \ref{Proposition 4.12}.
\end{proposition}
\vskip1mm
\begin{proof}
We have the map $\Phi_A: G_A\rightarrow \text{GL}(J/J^2)$, where the kernel is a unipotent group by Proposition \ref{Proposition 4.12}. Let $x\in \text{Im}(\Phi_A)$, then $\Phi_{A}^{-1}(x)=y.\text{Ker}(\Phi_A)$, where $\Phi_A(y)=x$. However, $\text{Ker}(\Phi_A)$ is unipotent, which is rational as a variety. Therefore, each fiber is a rational variety. Hence, $G_A$ and $\text{Im}(\Phi_A)$ are stably birationally equivalent. Moreover, using the Lemma \ref{Lemma 4.18}, we can conclude that $G_A\cong_b G_A/U\times U$, where $U$ is $\text{Ker}\Phi_A$. Hence, $G_A\cong_b \text{Im}(\Phi_A)\times U$, and $G_A(F)/R\cong \mathrm{Im}(\Phi_A)(F)/R$ for every field extension of $F/k$.
\end{proof}

\vskip2mm 
\noindent
Therefore, it is enough to study the image of the canonical map $\Phi_A$ for studying the rationality of $G_A$ for a split local algebra $A$.

\section{Finite Dimensional Commutative Algebras}\label{sec: commutative algebra}

\vskip2mm
\noindent
From this point onward, our primary emphasis will be on the automorphism group of split finite-dimensional local commutative associative algebras. Using the representation of finite-dimensional associative algebras by quiver and relation, one knows that any such algebra $A$ is isomorphic to $k[X_1,\dots, X_n]/{I}$, where $I$ is an ideal of $k[X_1,\dots, X_n]$ such that there exists an $l\geq2$ with $\langle X_1,\dots, X_n \rangle^l \subseteq I \subseteq \langle X_1,\dots, X_n \rangle^2$ (see \cite{AS1}, \cite{PG}). The minimum $l$ is called the Lowey length of $A$, denoted by $L(A)$. Every such ideal is called an \emph{admissible} or \emph{adequate ideal} for $A$ in $k[X_1,\dots,X_n]$. A finite set of generators for an adequate ideal is called \emph{adequate sets of relations} for $A$. When $A$ is realized as $k[X_1,\dots,X_n]/I$, then $J^m=(\text{Jacobson}(A))^m$ gets identified with $[(X_1,\dots,X_n)^m+I]/I$ and $n = \text{dim}(J/J^2)$. Let $A$ be identified with $k[X_1,\dots,X_n]/I$, where $I$ is an admissible ideal. Suppose we denote $x_i=X_i+I$, then $\{\bar{x_{i}}=x_i+J^2:i\in\{1,\dots,n\}\}$ is a basis of $J/J^2$. By taking matrices with respect to that basis, $\text{GL}(J/J^2)$ gets identified with $\text{GL}_n(k)$. The notion of change of variables is defined as follows.

\vskip1mm

\begin{definition}\label{Definition 5.1}$($\cite{AS1}$)$
A homomorphism of algebras $F:k[X_1,\dots,X_n]\longrightarrow k[X_1,\dots,X_n]$ is said to be a change of variables in $k[X_1,\dots,X_n]$, if there is a $n\times n$ matrix $(\lambda_{ij})\in \text{GL}_{n}(k)$ such that $F(X_j)=\sum_{1\leq i\leq n}\lambda_{ij}X_i$ modulo $(\langle X_1,\dots,X_n\rangle^2)$ for every $j\in \{1,\dots,n\}$. In particular, if no monomial of degree $\geq 2$ appears in the change of variables $F$, i.e., $F(X_j)=\sum_{1\leq i\leq n}\lambda_{ij}X_i$ for all $j\in\{1,\dots,n\}$, then we call such a change of variable a \emph{linear change of variables}, and $(\lambda_{ij})_{n\times n}$ is called the \emph{linear change of variables matrix} and denoted by $M_{F}=(\lambda_{ij})_{n\times n}$. For any linear change of variables $F$, we have $F(f)(X)=f(XM_{F})$, where $X=(X_1,\dots,X_n)$.
\end{definition}

\vskip2mm
\noindent
By using Theorem $1.5$ from \cite{Saorin} for a split commutative local algebra $A$, we get that $A$ is graded by the radical algebra if and only if $A$ admits a homogeneous admissible ideal $I$.


\vskip4mm

\begin{proposition}\label{Proposition 5.2}$($\cite{AS1}$)$  
Let A be a split commutative local algebra with Lowey length $l$ and $\emph{dim}(J/J^2) = n$. Suppose that $I, I'$ are two ideals of $k[X_1,\dots,X_n]$, the first being adequate for $A$. The following assertions are equivalent for a homomorphism of algebras
\begin{equation*}\label{eq19}
\phi: k[X_1,\dots,X_n]/I \longrightarrow k[X_1,\dots,X_n]/I':
\end{equation*}
\begin{enumerate}
\item $\phi$ is an isomorphism;
\vskip2mm
\item there is a change of variables $F$ in $k[X_1,\dots,X_n]$ such that $I'= F(I) + \langle X_1,\dots,X_n\rangle^l$ and $\phi(f+ I) = F(f) + I'$ for every $f\in k[X_1,\dots,X_n]$.  
\end{enumerate}
\vskip2mm
In particular, when $I = I'$, we get that every automorphism of $A$ is induced by a change of variables $F$ such that $F(I)\subset I$.
\end{proposition}
\vskip2mm
\noindent
We will denote by $C_{(A, I)}$ (resp. $CL_{(A, I)}$), the set of changes of variables (resp. linear changes of variables) $F$ of $k[X_1,\dots,X_n]$ satisfying $F(I)\subset I$. One knows that $C_{(A,I)}$ is a monoid and $CL_{(A,I)}$ is isomorphic to a subgroup $L_A$ of $G_A$ such that $L_A\cap \text{Ker}(\Phi_A)=\{1\}$ 
 (see \cite{AS3}, Corollary 21).
\vskip2mm
\noindent
Next, we review Lemma $22$ from \cite{AS3} in the context of a split local commutative algebra $A$ over $k$.

\vskip2mm

\begin{lemma}\label{Lemma 5.3}
$($\cite{AS3}, \emph{Lemma} $22)$ Let $A$ be a split local commutative algebra over $k$. Suppose $A\cong k[X_1,\dots,X_n]/I$ be a quiver representation with $I$ an \emph{admissible ideal} (a finite set of adequate relations). If $I$ is a homogeneous ideal (i.e., $A$ is a graded by the radical algebra), i.e, if homogeneous elements generate the \emph{adequate ideal} $I$, then
\begin{equation*}\label{eq20}
\emph{\text{ Im}}(\Phi_A )\cong CL_{(A, I)}\cong L_A.
\end{equation*}
\end{lemma}

\vskip2mm
\begin{proposition}\label{Proposition 5.4}
Let $A$ be a split local commutative graded by the radical algebra with a homogeneous admissible ideal $I$ and $\emph{dim}(J/J^2)=n$. Then, \[\emph{\text{Aut}}_k(A)\cong \emph{\text{Im}}(\Phi_A)\ltimes \emph{\text{Ker}}(\Phi_A), \text{ where $\Phi_A$ is the canonical map as in Proposition \ref{Proposition 4.12}. }\]
\end{proposition}
\vskip1mm
\begin{proof}
We have the canonical homomorphism $\Phi_A: \text{Aut}_k(A)\rightarrow \text{GL}(J/J^2)$ whose kernel is unipotent. Therefore, we have the following short exact sequence: 
\begin{center}
\begin{tikzcd}
1\arrow[r] & \text{Ker}(\Phi_A)\arrow[r] & \text{Aut}_k(A)\arrow[r] & \text{Im}(\Phi_A)\arrow[bend left=33]{l}{s}\arrow[r] & 1
\end{tikzcd}
\end{center}
It is known from Lemma \ref{Lemma 5.3} that any element of $\text{Im}(\Phi_A)$ gives an automorphism of $A$, which is induced by a linear change of variables. Hence, we have a section $s: \text{Im}(\Phi_A)\rightarrow \text{Aut}_k(A)$ given by $s(M)=\phi_M$, where $\phi_M(f+I)=f(XM)+I$, $X=(X_1,\dots,X_n)$. Finally, we have $\Phi_A\circ s=Id$, which further implies that the above sequence is split and provides the required isomorphism.
\end{proof}

\vskip2mm
\begin{proposition}\label{Proposition 5.5}
Let $A$ be a split local commutative graded by the radical algebra over $k$ with $\emph{dim}(J/J^2)=n$, then $G_A$ is $k$-isotropic, i.e., there exists a $k$-embedding of $\mathbb{G}_m$ in $G_A$. This $\mathbb{G}_m$ is a central torus in $\emph{\text{Im}}(\Phi_A)$ and contained in all maximal tori of $G_A$. In particular, $G_A$ can not be unipotent in this case.
\end{proposition}
\vskip1mm
\begin{proof}
Let $A$ be graded by the radical algebra, which implies that there exists a homogeneous adequate ideal, $I$ of $A$. We know that any automorphism $\phi$ of $A$ is induced by a change of variables $F$ such that $\phi(f+I)=F(f)+I$ and $F(I)\subseteq I$. Consider the embedding $\Gamma :\mathbb{G}_m\rightarrow G_A$ given by $\Gamma(\alpha)=\phi_{\alpha}$, where $\phi_{\alpha}(\bar X_i)=\alpha \bar X_i$, $\bar{X_i}\coloneq X_i+I$. Therefore, for each $\phi_{\alpha}$, the corresponding change of variables is $F_{\alpha}(X_i)=\alpha X_i$, which stabilizes $I$. This is an injective homomorphism. If $\alpha\in \mathbb{G}_m$ is in $\text{Ker}(\Gamma)$, then $(\alpha-1)X_i\in I\subseteq \langle X_1,\dots,X_n\rangle^2$ $\forall i$, which is only possible if $\alpha=1$. Hence, the group is $k$-isotropic. Moreover $\mathbb{G}_m$ embeds in $\text{Im}(\Phi_A)$ as $\text{Ker}(\Phi_A)$ is unipotent.
\vskip2mm
\noindent
Let $\phi_1\in G_A$ be such that it is given by a linear change of variables, i.e., $\phi(f+I)=F_1(f)+I$, where $F_1$ is a linear change of variables corresponding to $\phi_1$, and let $\phi_2\in \mathbb{G}_m$ with $F_2(X_i)=\alpha X_i$ as the corresponding change of variables. Then, it is easy to verify that $F_1\circ F_2=F_2\circ F_1$, as $F_1(f)(X)=f(\sum_{i=1}^n\lambda_{i1}X_i,\dots,\sum_{i=1}^{n}\lambda_{in}X_i)$, where $(\lambda_{ij})_{n\times n}$ is a non-singular matrix. It follows from Lemma \ref{Lemma 5.3} that $\text{Im}(\Phi_A)$ is given by a linear change of variables. Hence, via the above embedding, $\mathbb{G}_m$ is a central torus in $\text{Im}(\Phi_A$). Moreover, for any maximal torus $T$ in $G_A$, $T.\mathbb{G}_m$ is again a torus in $G_A$ containing $T$. Therefore, $\mathbb{G}_m\subset T$.  
\end{proof}


\vspace{0.2in}
\noindent
In the rest of the article, we will investigate the following questions for $G_A$, where $A$ is a split local finite-dimensional commutative associative algebra.

 
\vskip2mm

\begin{question}\label{Question 5.6}
Does $G_A$ always split, that is, whether the group always contains a split maximal torus?
\end{question}

\vskip2mm

\begin{question}\label{Question 5.7}
What is the necessary and sufficient condition on $A$ so that $G_A$ contains a split torus of rank strictly less than $\text{dim}(J/J^2)$?
\end{question}

\vskip2mm

\begin{question}\label{Question 5.8}
Is the group $G_A$ always $R$-trivial, when $A$ is a split local graded by the radical algebra?
\end{question}


\vskip2mm
\noindent
An affirmative answer to Question \ref{Question 5.6} implies a positive answer to Question \ref{Question 5.8}. However, the converse is not true. If $A$ is a local algebra with $\text{dim}(J/J^2) =1$, then the rank of $G_A$ is at most $1$ by Proposition \ref{Proposition 4.12}. So, if its rank is $1$ then by Proposition \ref{Proposition 5.5}, $G_A$ has to be split. We will show that Question \ref{Question 5.6} has no positive answer in general. This motivates us to answer Question \ref{Question 5.7}.

\vskip2mm

\begin{lemma}\label{Lemma 5.9}
Let $A$ be a split local commutative algebra over $k$ with $\emph{dim}(J/J^2)=n$. By quiver and relations representation of $A$ we have an isomorphism $A\cong k[X_1,\dots,X_n]/I$, where $\langle X_1,\dots,X_n \rangle^l \subseteq I\subseteq \langle X_1,\dots,X_n \rangle^2$, $l\geq 2$, $I$ is the admissible ideal of $A$. Then, $I=\langle X_1,\dots,X_n \rangle^l+\langle P_1,\dots,P_m \rangle$, where $l$ is the Lowey length of $A$ and all non-zero $P_i$'s are polynomials of degree at least $2$ and at most $l-1$ with no linear component.
\end{lemma}
\vskip1mm
\begin{proof}
Let $f\in I$. Then $f=\sum_{i_1,\dots,i_n}a_{i_{1},\dots,i_{n}}X_{1}^{i_{1}}\dots X_{n}^{i_{n}}$. By the given hypothesis, we have $\langle X_1,\dots, X_n \rangle^l\subseteq I$; therefore, monomials with degree greater than or equal to $l$ are contained in $I$. Hence, we can write $f=\sum_{i}\lambda_{i}h_i+P$, where $h_i$'s are distinct monomials in \{$X_1,\dots,X_n$\} of degree at least $l$ and $P$ is a polynomial of degree at most $l-1$. However, $f-\sum_{i}\lambda_{i}h_{i}\in I$, therefore $P\in I$. Since $I$ is contained in $\langle X_1,\dots,X_n\rangle ^2$, $P$ must have degree at least $2$ with no linear component. Now, for each $f\in I$, we have a unique $P$ in $I$ which has the degree at least $2$ and at most $l-1$ with no linear component. Consider the set
\[S=\{P: P \text{ corresponds to each } f\in I \text{ with the above property }\}.\] Suppose $I$ is a proper ideal lying in between $\langle X_1,\dots,X_n\rangle^l$ and $\langle X_1,\dots,X_n\rangle ^2$. Then $S$ must be non-empty. Therefore, the ideal generated by $S$ is contained in $I$, and let us call it $I_1$. Moreover, $k[X_1,\dots,X_n]$ is a noetherian ring, so $I_1= \langle P_1,\dots,P_m \rangle$ where $P_i\in S$ and $I_1\subset I$. Hence $I=\langle X_1,\dots,X_n\rangle^l+\langle P_1,\dots,P_m\rangle$, where ${P_i}$'s have the property given in the hypothesis.
\end{proof}
\vskip2mm
\begin{definition}\label{Definition 5.10} 
Stabilizer subgroup of a polynomial $f\in k[X_1,\dots,X_n]$ in $\text{GL}(J/J^2)$, where $\text{dim}(J/J^2)=n$, is defined by 
\begin{equation*}\label{eq28}
\text{Stab}(f)\coloneq\{M\in \text{GL}(J/J^2): f(XM)=f(X), X=(X_1,\dots,X_n)\}.
\end{equation*}
\end{definition}
\noindent
$\text{Stab}^{0}(f)$ denotes the identity component of the stabilizer subgroup of $\text{GL}(J/J^2)$.

\vskip4mm
\begin{definition}\label{Definition 5.11}
The similarity subgroup of a polynomial $f\in k[X_1,\dots,X_n]$ in $\text{GL}(J/J^2)$, where $\text{dim}(J/J^2)=n$, is defined by \[\text{Sim}(f)\coloneq\{M\in \text{GL}(J/J^2): f(XM)=\alpha_M f(X), X=(X_1,\dots,X_n)\text{ and } \alpha_M\in k^{\times}\}.\]
\end{definition}
\noindent
$\text{Sim}^{0}(f)$ denotes the identity component of the similarity subgroup in $\text{GL}(J/J^2)$. It is easy to see that $\text{Sim}^{0}(f)=\mathbb{G}_m.\text{Stab}^{0}(f)$.


\vskip2mm
\begin{proposition}\label{Proposition 5.12}
Let $A$ be a split local commutative graded by the radical algebra with $\emph{dim}(J/J^2)=n$. Suppose the corresponding admissible ideal for $A$ is $I\cong\langle X_1,\dots, X_n \rangle^l+\langle f\rangle$, where $2\leq \emph{\text{deg}}(f)< l$ is the Lowey length of $A$ and $f$ is a homogeneous polynomial.  Then,
\begin{equation*}\label{eq23}
\emph{Aut}_k(A)\cong \emph{Sim}(f)\ltimes \emph{\text{Ker}}(\Phi_A), where \text{ $\Phi_A$ is the canonical map as in Proposition \ref{Proposition 4.12}. }
\end{equation*}
\end{proposition}

\vskip1mm
\begin{proof}
We have the canonical map $\Phi_A: \text{Aut}_k(A)\rightarrow\text{GL}(J/J^2)$. It is enough to prove that for the given $I$, $\text{Im}(\Phi_A)=\text{Sim}(f)$. Let $M\in \text{Im}(\Phi_A)$, then $M$ induces an automorphism of  $A$ which is given by the linear change of variables $F_1(X)=XM$, where $X=(X_1,\dots,X_n)$ and $F_1(I)\subset I$ by Lemma \ref{Lemma 5.3}. As $f\in I$, \[F_1(f)(X)=f(\sum_{i}m_{i1}X_i,\dots,\sum_{i}m_{in}X_i)\in I,\] where $M=(m_{ij})_{n\times n}$ and $X=(X_1,\dots,X_n)$. Since a linear change of variables does not change the degree of the polynomial, we have $\text{deg}(F(f)=\text{deg}(f)$ and $F_1(f)$ is a homogeneous polynomial. As $f$ is a non-zero homogeneous polynomial, we have $F_1(f)=gf$, for some non-zero homogeneous polynomial $g\in k[X_1,\dots,X_n]$. However, $\text{deg}(f)=\text{deg}(F_1(f))$ implies $\text{deg}(g)=0$. Therefore,
$F_1(f)=\alpha f$ and $f(\sum_{i}m_{i1}X_i,\dots,\sum_{i}m_{in}X_i)=\alpha f(X_1,\dots,X_n)$ with $\alpha\in k^{\times}$.
Hence, $M\in \text{Sim}(f)$. So, we have $\text{Im}(\Phi_A)\subset \text{Sim}(f)$. It is easy to see that $\text{Sim}(f)\subset \text{Im}(\Phi_A)$. Now using Proposition \ref{Proposition 5.4}, $\text{Aut}_k(A)\cong \text{Sim}(f)\ltimes \text{Ker}(\Phi_A)$.
\end{proof}

\vskip2mm

\begin{theorem}\label{Theorem 5.13}
Let $A$ be a finite-dimensional local commutative algebra over a field $k$ with $\text{dim}(J/J^2)=n$. Suppose $A\cong k[X_1,\dots, X_n]/I$ by quiver and relation representation of $A$, where $I=\langle X_1,\dots, X_n\rangle^l+\langle f\rangle$ and $2\leq \mathrm{ deg}(f)<l$ is the Lowey length of $A$. Then, $G_A$ is rational or $R$-trivial as a $k$-group if and only if $\mathrm{Stab}^{0}(f)$ is rational or $R$-trivial.
\end{theorem}
\vskip1mm
\begin{proof}
Applying Proposition \ref{Proposition 5.12}, we get that $G_A\cong \mathbb{G}_m.\text{Stab}^{0}(f)\ltimes \text{Ker}(\Phi_A)^{0}$, where $\text{Ker}(\Phi_A)^{0}$ is a unipotent group. Therefore, $G_A$ is either rational or $R$-trivial as $\text{Stab}^{0}(f)$ is either rational or $R$-trivial.
\end{proof}

\vskip2mm
\begin{proposition}\label{Proposition 5.14}
$($\cite{CM}, \emph{Proposition} $2.4)$ Unitary, special orthogonal groups, and symplectic groups are rational.   
\end{proposition}

\vskip2mm
\begin{example}\label{Example 5.15}
Let $I$ be an admissible ideal of $k[X_1,\dots,X_n]$, which equals to $\langle X_1,\dots, X_n \rangle^l+\langle q\rangle$, where $q$ is a non-degenerate quadratic form, $\text{char}(k)\neq2$ and $l>2$. Then $G_A$ is rational for $A\cong k[X_1,\dots,X_n]/I$ by using Theorem \ref{Theorem 5.13} as $\text{Stab}^{0}(q)\cong\text{SO}(n,q)$ which is a simple rational $k$-group by Proposition \ref{Proposition 5.14}. Similarly, if $G=\mathrm{Sp}_n$, then also we can construct an algebra $A$ such that $G_A$ is $R$-trivial. These examples show that the conditions in Theorem \ref{Theorem 4.20} are sufficient but not necessary.
\end{example}

\vskip2mm 
\noindent
Let $(V,q)$ be a quadratic space. Recall that $\text{SO}(V,q)$ is a $k$-isotropic group if and only if $q$ is isotropic over $k$ (\cite{Borel}, Section $23.4$).

\vskip4mm
\begin{theorem}\label{Theorem 5.16}
Let $A$ be a finite-dimensional local commutative algebra with $\text{dim}(J/J^2)=n$. Suppose $A\cong k[X_1,\dots, X_n]/I$ by quiver and relation representation of $A$, where $I$ is given by $\langle X_1, X_2,\dots, X_n\rangle^l+\langle q\rangle$ and $\emph{char }k\neq2$; $l>2$ is the Lowey length of $A$ and $(J/J^2,q)$ is an anisotropic quadratic space over the field $k$, then $G_A$ is not a $k$-split group.
\end{theorem}
\vskip1mm
\begin{proof}
We know from Proposition \ref{Proposition 5.12}, that $\text{Aut}_k(A)\cong \text{Sim}(q)\ltimes \text{Ker}(\Phi_A)$, where $\text{Ker}(\Phi_A)$ is unipotent, and \[\text{Sim}(q)=\{M\in \text{GL}(J/J^2):q(XM)=\alpha_M q(X), X=(X_1,\dots, X_n)\text{ and }\alpha_M\in k^{\times}\}.\] However, we have the following short exact sequence;
\begin{center}
\begin{tikzcd}
1\arrow[r] & \text{Stab}(q)\arrow[r] & \text{Sim}(q)\arrow[r,"\beta"] & \mathbb{G}_{m,k}\arrow[r] & 1
\end{tikzcd} with $\beta(M)=\alpha_M$.
\end{center}
Therefore, $\text{Sim}(q)= \text{Stab}(q).\mathbb{G}_m\cong \text{SO}(n,q).\mathbb{G}_m$. This implies that any maximal torus of $G_{A}$ is isomorphic to $S.\mathbb{G}_m$, where $S$ is a torus in $\text{SO}(n,q)$. However, $S$ can not be isotropic because if $S$ is $k$-isotropic, then $\text{SO}(n,q)$ is $k$-isotropic, which implies $q$ is isotropic over $k$, which is a contradiction. Hence, there is no split maximal torus in $G_{A}$.
\end{proof}

\vskip2mm

\begin{remark}\label{Remark 5.17}
This result shows that Question \ref{Question 5.6} is not true in general unless the $\text{dim}(J/J^2)$ is $1$.    
\end{remark}

\vskip2mm
\noindent
We can produce more examples of non-split $G_A$ using the following proposition.

\vskip4mm
\noindent
Let $A$ be an arbitrary basic not necessarily commutative algebra $A$. If $\Gamma$ is the quiver of $A$ and $I$ is an ideal of $k[\Gamma]$, where $k[\Gamma]$ is the path algebra generated by all paths and $I_{*}$ is the ideal of $k[\Gamma]$ generated by the homogeneous components with the smallest degree of elements of $I$ (see \cite{AS1}, \cite{AS2}). We would like to state the following proposition to use in the future.

\vskip2mm

\begin{proposition}\label{Proposition 5.18}
$($\cite{AS1}, \emph{Proposition} $2.1)$ Let $A$ be a basic not necessarily commutative algebra, with quiver $\Gamma$ and $G_{r}(A)$ be the associated graded by the radical algebra of $A$. If $I$ is an adequate ideal for $A$ in $k[\Gamma]$, then $I_{*}$ is an adequate ideal for $G_{r}(A)$. Conversely, if $L$ is a homogeneous adequate ideal for $G_{r}(A)$, then there is an adequate ideal $I$ for $A$ such that $I_{*}=L$.
\end{proposition}

\vskip2mm
\begin{corollary}\label{Corollary 5.19}
Let $A$ be a split local commutative algebra with the given admissible ideal of the form $I\cong \langle X_1,\dots, X_n\rangle^l+\langle f_1,\dots,f_m\rangle $, where $f_i=\sum_{i_1,i_2,...i_n} a_{i_1,...i_n}X_{1}^{i_{1}}\dots X_{n}^{i_{n}}+b_iq$; $b_i, a_{i_1,...i_n}\in k^{\times}$ and $\emph{deg}(f_i)<l$, and $q$ is a fixed anisotropic quadratic form. Then $k$-split rank of $\emph{\text{Aut}}_k(A)$ is at most $1$.
\end{corollary}
\vskip1mm
\begin{proof}
Proposition \ref{Proposition 5.18} shows that the adequate ideal for the associated graded by the radical algebra of $A$ is $I_{*}=\langle X_1, X_2,\dots,X_n \rangle^l+\langle q\rangle$. Therefore, $G_{r}(A)\cong k[X_1,\dots, X_n]/I_{*}$. It follows from Theorem \ref{Theorem 5.16} that any maximal torus in $\text{Aut}_k(G_r(A))$ is isomorphic $\mathbb{G}_m.S$, where $S$ is a maximal torus in $\text{SO}(n,q)$. Hence, the $k$-split rank of $\text{Aut}_k(G_{r}(A))$ is $1$. As in Proposition \ref{Proposition 4.12}, we consider the map $\Psi_A: \text{Aut}_k(A) \rightarrow \text{Aut}_k(G_{r}(A))$, which possesses a kernel that is unipotent. Hence, any torus of $\text{Aut}_k(A)$ embeds in $\text{Aut}_k(G_{r}(A))$, therefore, the $k$-split rank of $\text{Aut}_k(A)$ is at most $1$.
\end{proof}
\vskip2mm
\noindent 
Now we define some terminology before proceed to our next result.
\vskip2mm
\begin{definition}\label{Definition 5.20}
A non-constant homogeneous polynomial that is not a monomial is said to be an \emph{$s$-homogeneous polynomial} if it is generated by indeterminants $X_s,\dots, X_n$, $s\geq 1$. 
\end{definition}
\vskip4mm
\begin{example}\label{Example 5.21}
Let $g=X_3^8+X_4^8\in k[X_1,\dots, X_4]$ be a $3$-homogeneous polynomial.
\end{example}
\vskip4mm
\begin{definition}\label{Definition 5.22}
A non-constant homogeneous polynomial $f$ is said to be \emph{$s$-monomial homogeneous polynomial} if $f=Mg$, where $M$ is a monomial or a scalar and $g$ is an $s$-homogeneous polynomial for $s\geq 1$.
\end{definition}
\vskip4mm
\begin{example}\label{Example 5.23}
Let $f=X_1^2X_2^3X_3^4 X_4^8+X_1^2X_2^3X_3^{12}\in k[X_1,\dots,X_4]$. Then it is a $3$-monomial homogeneous polynomial because $f=X_1^2X_2^3X_3^4(X_3^8+X_4^8)$.
\end{example}
\vskip4mm
\begin{definition}\label{Definition 5.24}
Let $A$ be a split local commutative algebra with Lowey length $l$ and $\text{dim}(J/J^2)=n$ over $k$. It has been proved in Lemma \ref{Lemma 5.9} that the corresponding admissible ideal $I$ for $A$ is generated by all monomials of degree $l$ and non-zero polynomials $P_i$ of degree at least $2$ and at most $l-1$ with no linear component, i.e., $I=\langle X_1,\dots,X_n\rangle^l+\langle P_1,\dots, P_m\rangle$ for some $m$, where $2\leq \text{deg}(P_i)\leq l-1$ and $P_i$ has no linear component for all $i$. If all $P_i$'s are monomials or $0$, then we will say that $I$ is a monomial ideal. The admissible ideal $I$ satisfies the \emph{Property $*$ for $r$}, where $r\in \mathbb{N}$, if there exists a $P_i$ which is an $s_i$-monomial homogeneous polynomial for some $s_i\geq r$, and $P_j$ is either a monomial or an $u_j$-monomial homogeneous polynomial for some $ u_j\geq r$ for all $j\neq i$.
\end{definition}

\vskip4mm

\begin{theorem}\label{Theorem 5.25}
Let $A$ be a split local commutative algebra over an infinite field $k$ with $\emph{dim}(J/J^2)=n$. Let $r\in \mathbb{N}$ and $1\leq r\leq n$. Then the following statements are equivalent:
\begin{enumerate}
    \vskip4mm
    \item $G_A$ contains a diagonal change of variables of the form \[D(r,k)=
\begin{pmatrix}
a_1 & 0 &\cdots & 0 & 0 & \cdots & 0 \\
0 & a_2 & \cdots & 0 & 0 & \cdots & 0\\
\vdots & \vdots & \ddots & \vdots& \vdots & \vdots & \vdots\\
0 & 0 & \cdots & a_{r-1} & 0 &\cdots & 0 \\
0 & \cdots & 0 & 0 & \beta & \cdots & 0\\
\vdots & \vdots & \vdots & \vdots & \vdots & \cdots & \vdots\\
0 & \cdots & \cdots & 0 & 0 & \cdots & \beta
\end{pmatrix}_{n\times n}
,\]
where $a_i,\beta\in k^{\times}, 1\leq i\leq r-1;$
\vskip4mm
     \item There exists an admissible ideal $I$ for $A$ which satisfies the Property $*$ for $r$ or a monomial ideal $I$ such that $A\cong k[X_1,\dots,X_n]/I$.
\end{enumerate}
\vskip4mm
If the hypothesis $(2)$ is satisfied, we have a $k$-embedding $\phi: \mathbb{G}_m^{r}\rightarrow G_A$, where $r\leq n$. 
\end{theorem}
\vskip1mm
\begin{proof}
$\bigl(\bold{(2)}\Rightarrow \bold{(1)}\big)$
Let us denote $X_i+I$ by $\bar{X_i}$. Under the hypothesis $(2)$, consider the following map \[\phi:\mathbb{G}_m\times \mathbb{G}_m\dots \times \mathbb{G}_m(r \text{ times }) \rightarrow G_A\] which is given by
\begin{equation}\label{eq 4.1.9}
    \phi(a_1,\dots,a_{r-1},\beta)=
    \begin{cases}
         F(\bar{X_i})=a_i\bar{X_i} \text{ if } 1\leq i \leq r-1 \\
     F(\bar{X_i})=\beta \bar{X_i} \text{ if } r\leq i\leq n
    \end{cases}
\end{equation}    
It is given that $I=\langle X_1,\dots,X_n\rangle^l+\langle P_1,\dots,P_m\rangle$ for some $m$ and $l=\text{ Lowey length of }A$, where $P_i$ is either a monomial or an $s_i$-monomial homogeneous polynomial of degree at least $2$ and atmost $l-1$, $s_i\geq r\geq1$, for all $1\leq i\leq m$. Therefore, the corresponding change of variables $F$ stabilises $I$, for all $F\in D(r,k)$. Hence, we have $\phi(a_1,\dots a_{r-1},\beta)\in G_A$ and $D(r,k)\subset G_A$. Suppose $\phi(a_1,\dots,a_{r-1},\beta)=\text{Id}$, then \[(a_i-1)X_i\in I \quad \text{if } 1\leq i\leq r-1, \text{ and } (\beta-1)X_i\in I \text{ if } r\leq i \leq n\] Since $I\subset \langle X_1,\dots,X_n\rangle^2$, we have $\text{Ker}(\phi)$ is trivial. Hence, the above map (\ref{eq 4.1.9}) is a $k$-embedding.
\vskip6mm
\noindent
$\big(\bold{(1)}\Rightarrow \bold{(2)}\big)$ 
\noindent
(\textbf{Brief Idea of the Proof:}) So to prove this part, we use the induction hypothesis on the Lowey length of $A$, $L(A)$. Given that $G_A$ satisfies a hypothesis, we will show that $A$ satisfies some property. For that, we construct a new algebra $\tilde{A}$ from $A$ which has Lowey length $L(A)-1$ and $G_{\tilde{A}}$ satisfies the same hypothesis as $G_A$. Therefore, by the induction method, $\tilde{A}$ satisfies the same property as $A$. Then, using the property of $\tilde{A}$, we prove that the same property is true for $A$.
\vskip2mm
\noindent
(\textbf{Induction Statement:}) Assume $G_A$ contains a class of automorphisms which are given by the change of variables \[\biggl\{D(r,k)=\{\Gamma(\bar{X_i})=a_{i}\bar{X_{i}} \text{ if } 1\leq i\leq r-1 \text{ and } \Gamma(\bar{X_i})=\beta\bar{X_i} \text{ if } r\leq i\leq n\}\biggl\}.\] Then we prove that the admissible ideal $I$ for $A$ satisfies the Property $*$ for some $r$ or $I$ is a monomial ideal (given in Definition \ref{Definition 5.24}). To prove this, we use induction on $l$, where $l$ is the Lowey length of $A$. Lowey length of $A$ is given by $l$=min $\{k\geq 2:\langle X_1,\dots,X_n \rangle^k \subseteq I \}$.\par
\vskip6mm
\noindent
(\textbf{True for Lowey length $l=2$}:) If $l=2$, then $I=\langle X_1,\dots, X_n\rangle^2$, which is a monomial ideal, and $D(r,k)\subset G_{A}\cong \text{GL}_n$. Now let us assume that $I$ is properly contained in between $\langle X_1,\dots,X_n\rangle^l$ and $\langle X_1,\dots,X_n \rangle^2$ and $l>2$. Let $J(A)$ be the Jacobson radical of $A$ and we have $(J(A))^m=(\langle X_1,\dots,X_n\rangle^m+I)/I$ for all $m\geq 1$. Set $\tilde{I}=I+\langle X_1,\dots,X_n\rangle^{l-1}$ and let $\Tilde{A}=A/J(A)^{l-1}$. Then, $\Tilde{A}$ has lowey length $l-1$ with $\tilde{I}$ is an admissible ideal for $\Tilde{A}$. 
\vskip4mm
\noindent
(\textbf{True for Lowey length $l-1$}:) Suppose, the rank of $G_{\Tilde{A}}=t$. We know that there is a homomorphism from $G_A\rightarrow G_{\Tilde{A}}$ whose kernel is unipotent by Proposition \ref{Proposition 4.12}. Hence, we have an embedding of change of variables $D(r,k) \subset G_{\Tilde{A}}$ and $r\leq t\leq n$. Therefore by the induction hypothesis, $\tilde{I}$ satisfies Property $*$ for $r$ or $\tilde{I}$ is a monomial ideal for $\Tilde{A}$, as $\Tilde{A}$ has Lowey length $l-1$ and $G_{\tilde{A}}$ contains $D(r,k)$. Hence, $\tilde{I}$ is generated by all monomials of degree $l-1$ and $S_i$, where $S_i$ is either an $u_i$-monomial homogeneous polynomial or a monomial, $u_i\geq r$, for all $1\leq i\leq m$, and $2\leq$ $\text{deg}(S_i)\leq l-2$, i.e., $\tilde{I}=\langle X_1,\dots,X_n\rangle^{l-1}+\langle S_1,\dots,S_m\rangle$.
\vskip4mm
\noindent
(\textbf{Claim: The statement is true for Lowey length $l$}:) To prove $I$ satisfies Property $*$ for $r$ or $I$ is a monomial ideal, it is enough to prove that every generator of $I$ of degree at least $2$ and at most $l-1$ is either an $s$-monomial homogeneous polynomial for $s\geq r$ or a monomial. That is if $I=\langle X_1,\dots, X_n\rangle^l+\langle f_1,\dots,f_m\rangle$, where $2\leq \text{deg}(f_i)\leq l-1$ and $f_i$ has no linear part, then $f_i$ is either $s_i$-monomial homogeneous polynomial for $s_i\geq r$ or a monomial. 
\vskip2mm
\noindent
If $f_i$ is either a monomial or an $s_i$-monomial homogeneous polynomial for $s\geq r\geq 1$, then we are done.  Suppose there exists a generator $f_j$ of $I$ such that $f_j=\sum_{i}\alpha_{i}h_{i}+ g$, where $h_{i}$'s are distinct monomials of degree $l-1$, $\alpha_i \in k$ and $2\leq\text{deg}(g)<l-1$. However, $f_j\equiv g \ (\text{mod }\langle X_1,\dots,X_n\rangle^{l-1})$. Hence, we have $g\in \tilde{I}$ is a generator in $\tilde{I}$ of degree at least $2$ and at most $l-2$, as $\tilde{I}=I+\langle X_1,\dots,X_n\rangle ^{l-1}$ and $f_j$ is a generator of $I$ with $2\leq \text{deg}(f_j)\leq l-1$. Since $\tilde{I}$ satisfies the Property $*$ for $r$ or $\tilde{I}$ is a monomial ideal, we have $g$ is either an $u$-monomial homogeneous polynomial for $u\geq r$ or a monomial. If $g$ is a $u$-monomial homogeneous polynomial, then let $g=MP$ with $M$ a monomial or a scalar, and $P$ is a $u$-monomial homogeneous polynomial of degree $d$.  However, $f_j=0$ in $A$, we have 
\begin{equation}\label{eq 5.2}
g(\bar{X_1},\dots,\bar{X_n})=-\sum_i\alpha_{i}h_{i}(\bar{X_1},\dots,\bar{X_n})
\end{equation}
in $A$. Moreover, $\Gamma(I)\subset I$ and apply $\Gamma$ on both sides of the previous equation (\ref{eq 5.2}). So, we have a function $\Delta_g$ depending on $g$ and $\Gamma$ defined below such that;
\begin{align}\label{eq 5.3}
\begin{split}
(\Delta_{g})g(\bar {X_1},\dots,\bar{X_n})=-\sum_i\alpha_ih_i(a_1,\dots,a_{r-1},\beta,\dots,\beta)h_i(\bar{X_1},\dots,\bar{X_n}).
\end{split}
\end{align}
where \begin{equation*}
    \Delta_{g}=
   \begin{cases}
        M(a_1,\dots,a_{r-1},\beta,\dots,\beta)\beta^d ,\text{ if } g=MP \text{ is a }$u-$\text{monomial homogeneous polynomial.}\\
        g(a_1,\dots,a_{r-1},\beta,\dots,\beta), \text{ if } g \text{ is monomial. }
    \end{cases}
\end{equation*}
\vskip2mm
\noindent
Again, by substituting $g(\bar{X_1},\dots,\bar{X_n})=-\sum_i\alpha_{i}h_i(\bar{X_1},\dots,\bar{X_n})$ in equation (\ref{eq 5.3}), we have 
\begin{align*}
\begin{split}
(\Delta_{g})\sum_i\alpha_ih_i(\bar{X_1},\dots,\bar{X_n})=\sum_i\alpha_ih_i(a_1,\dots,a_{r-1},\beta,\dots,\beta)h_i(\bar{X_1},\dots,\bar{X_n}),
\end{split}
\end{align*}
\noindent
Since ${h_i}$'s are linearly independent in $J(A)^{l-1}$, we get that
\begin{equation*}
(\Delta_g)\alpha_i=\alpha_ih_i(a_1,\dots,a_{r-1},\beta,\dots,\beta).
\end{equation*}
This is true for all $\{a_1,\dots,a_{r-1},\beta\}\in k$ and $\forall i$. As $k$ is infinite, $\alpha_{i}=0, \forall i$. Therefore, we get that every generator in $I$ of degree at least $2$ and at most $l-1$ is either an $s$-monomial homogeneous polynomial for $s\geq r$ or a monomial. Hence, we are done.
\end{proof}

\vskip1mm

\begin{remark}\label{Remark 5.26}
One can establish that any permutation of the specified change of variables, accompanied by different weights, will ensure that the admissible ideal $I$ retains the Property $*$ for some $r$ with the designated weights for the indeterminates. We can get Theorem $1.2$ of \cite{AS1} as a particular case of the above theorem.
\end{remark}
\vskip1mm

\begin{corollary}\label{Corollary 5.27}
Let $A$ be a split local commutative algebra with $I$ an admissible ideal and $\emph{dim}(J/J^2)=n$, then $I$ is monomial if and only if $\mathbb{G}_m^n=D(n,k)$ embeds in $G_A$. 
\end{corollary} 
\vskip1mm
\begin{proof}
    Take $D(n,k)=\text{diag }\{(a_1,\dots,a_n): a_i\in k^{\times}\}$, then follow the previous technique with $r=n$.
\end{proof}
\vspace{0.1in}
\noindent
So, to study the automorphism group of split local commutative associative algebras, we see that it is important to examine the stabilizer of polynomials, where these polynomials occur as the generator of the ideal $I$ of $A$ in quiver and relation representation of $A$. Therefore, we recall now an important result related to the stabilizer of a polynomial from the paper \cite{GG}.

\vspace{0.2in}

\begin{definition}\label{Definition 5.28} 
Stabilizer of a polynomial $f\in k[V]$ is defined by
\begin{equation*}\label{eq28}
\text{Stab}(f)=\{g\in \text{GL}(V): f\circ g=f\}
\end{equation*}
\end{definition}

\vskip2mm

\begin{lemma}\label{Lemma 5.29}
$($\cite{GG}, \emph{Lemma} $5.1)$ Let $X\subset \emph{SL}(V)$ be a simple algebraic group over an algebraically closed field $k$ such that $V$ is irreducible, restricted, and tensor indecomposable for $X$. Put $q$ for a non-zero $X$ invariant quadratic form on $V$ if one exists; otherwise, set $q\coloneq0$. If $(X,V)$ does not appear in Table $1$ of \cite{Sei}, then, for every $f\in k[V]^X-k[q]$, the stabilizer of $f$ in $\emph{\text{GL}}(V)$ has identity component $X$.
If additionally, $\emph{char}(k)=p\neq 2,3$ and does not divide $\emph{\text{deg}}(f)$,  $f$ is not in $k[V]^{(p)}[q]$ $($see \cite{GG} for notation$)$, and furthermore $\emph{char}(k)$ does not divide $n+1$ if $X$ has type $\emph{A}_n$, then the scheme-theoretic stabilizer of $f$ in $\emph{\text{GL}}(V)$ is smooth with identity component $X$. The above lemma is also true over an arbitrary base field.
\end{lemma}

\vskip2mm

\begin{proof}
Let $k$ be an arbitrary field. First, we choose a pair $(X=G, V)$ to satisfy the above hypothesis. The natural homomorphism \[k[V]^G\otimes k_{alg}\rightarrow k_{alg}[V\otimes k_{alg}]^{G\times k_{alg}}\] is an isomorphism. So, there exists an $f\in k[V]^G-k[q]$. Therefore, $G(k)=G(k_{alg})\cap \text{GL}(V)=\text{Stab}^{0}(f)(k_{alg})\cap \text{GL}(V)=\text{Stab}^{0}(f)(k).$
\end{proof}

\vskip2mm
\begin{definition}\label{Definition 5.30}
  $($\cite{OS}, \cite{Schneider}$)$  A non-zero homogeneous polynomial $f\in k[X_1,\dots, X_n]$ of degree $d$ is called non-singular if any one of the following equivalent conditions holds,
    \begin{enumerate}
        \item $0=(0,\dots,0)$ is the unique common root of its partial derivatives $\displaystyle\frac{\partial f}{\partial
        X_1},\dots,\displaystyle\frac{\partial f}{\partial X_n}$. 
        \vskip2mm
        \item Discriminant $(\Delta_{n,d}(f))$ of $f$ does not vanish (see Chapter $13$, \cite{GKZ}).
        \vskip2mm
        \item Let $\Theta_f$ be the corresponding symmetric multilinear form of $f$ of degree $d$ over a vector space $V=k^n$ of dimension $n$ over $k$. Then $\Theta_f(v,\dots,v,w)=0$ for all $w\in V$ implies $v=0$. 
    \end{enumerate}
\end{definition}

\vspace{0.1in}
\noindent
Next, we recall the following theorem, originally due to Camille Jordan.
\begin{theorem}\label{Theorem 5.31}
$($\cite{Schneider}$)$ Let $\Theta$ denote a symmetric multilinear map of degree $r\geq 3$ on a vector space of dimension $n$ over a field $k$. Assume that $k$ has characteristic $0$ or greater than $r$. If $\Theta$ is nonsingular, then its orthogonal group (stabilizer group) in $\emph{GL}_n(k)$ is finite.
\end{theorem}
\vskip2mm
\noindent
We now prove the following result.
\vskip 2mm
\begin{proposition}\label{Proposition 5.32}
Let $A$ be a split local commutative graded by the radical algebra with Lowey length $l$ and $\emph{dim}(J/J^2)=n$, where the admissible ideal $I$ is given by $\langle X_1,\dots,X_n \rangle^l +\langle f_1,\dots,f_m\rangle$, $f_i$'s are homogeneous polynomials such that not all $f_{i}$'s are zero. Let $d_i=\emph{deg}(f_i)$, where $1\leq i\leq m$ and $2\leq d_i<l$.  Suppose there exists a unique generator $f_s$ with degree $\emph{min}\{d_i: 1\leq i\leq m\}$. Then 
\[\bigcap_{i=1}^{m}\emph{Sim}(f_i)\subset\emph{\text{Im}}(\Phi_A)\subset \emph{Sim}(f_s).\]
In particular, if $f_s$ is a non-singular homogeneous polynomial of degree $\geq 3$ over a field $k$ with $\emph{char }k=0$ or $\emph{char }k=p>3$ and $f_s$ is non-singular polynomial of degree $d$ with $2<d<p$, then $G_A$ is a rational group of rank $1$.
\end{proposition}
\vskip1mm
\begin{proof}
 It is known from Lemma \ref{Lemma 5.3} that when the admissible ideal $I$ for $A$ is a homogeneous ideal, we have \[\text{Im}(\Phi_A)=\{M\in \text{GL}(J/J^2): \phi_M \in \text{Aut}_k(A) \text{ with } \phi_M(f+I)=f(XM)+I \}.\] Therefore, every $M\in\text{Im}(\Phi_A)$ induces an automorphism of $A$ which is given by the linear change of variables $F_1(X)=X_1M$, where $X=(X_1,\dots,X_n)$ and $F_1(I)\subset I$ by Proposition \ref{Proposition 5.2} and Lemma \ref{Lemma 5.3}. It is easy to see that $\bigcap_{i=1}^{m}\text{Sim}(f_i)\subset \text { Im}(\Phi_A)$. Let us denote $f_M(X)\coloneqq f(XM)$ and a linear change of variables by $F_1$ for $M$, where $M\in \text{ Im}(\Phi_A$). Since a linear change of variables does not change the degree of the polynomial and $F_1(f_i)={f_{i}}_{M}\in I$ by Lemma \ref{Lemma 5.3}, we have ${f_{i}}_M$ is a homogeneous polynomial of degree $d_i$ and ${f_{i}}_M=F_1(f_i)=\sum_{j=1}^{m}g_{ij}f_j$, where $g_{ij}\in k[X_{1},\dots,X_{n}]$ is a homogeneous polynomial of degree $s_{ij}$. As $F_1$ is given by the linear change of variables, we have
\begin{equation*}\label{eq30} 
d_i= \text{deg}({f_{i}}_{M})=\text{deg}(F_1(f_i))=\text{deg }(\sum_{j=1}^{m}g_{ij}f_j)=\text{ Max}_{j=1}^{m}(\text{deg}(g_{ij})+\text{deg}(f_j)).
 \end{equation*}
 Without loss of generality, let $d_1=\text{ min }d_i$, $d_i$ is degree of $f_i$ and $d_i
 \leq d_{i+1}$ for all $i\geq 2$. Since ${f_{i}}_M$ is a homogeneous polynomial of degree $d_i$, we have $d_i=(s_{ij}+d_j)$ for all $1\leq i,j\leq m$, where $\text{deg}(g_{ij})=s_{ij}$. By the given hypothesis, $f_1$ is the unique generator of minimum degree $d_1$. Therefore, $d_1=s_{11}+d_1$ and $d_1<d_i$ for all $i\neq 1$, which implies $s_{11}=0$ and $g_{1j}=0,\forall j\in\{2,\dots,m\}$. So, we get $f_1(XM)=\alpha f_1$. Hence, $\text{Im}(\Phi_A)\subset \text{Sim}(f_1)$. Finally, we have \[\bigcap_{i=1}^{m}\text{Sim}(f_i)\subset \text{Im}(\Phi_A)\subset \text{Sim}(f_1).\] Now from Proposition \ref{Proposition 5.4}, we know that $G_A\cong \text{Im}(\Phi_A)^{0}\ltimes \text{Ker}(\Phi_A)^{0}$. It is clear from Theorem \ref{Theorem 5.31} that under the assumption on $f_1$ depending of characteristic of $k$, $\text{Sim}^{0}(f_1)=\mathbb{G}_m$ as $\text{Stab}^{0}(f)=\{Id\}$. Hence, $\text{Im}(\Phi_A)^{0}\cong \mathbb{G}_m$. Therefore, $G_A\cong \mathbb{G}_m\ltimes \text{Ker}(\Phi_A)^{0}$, hence $G_A$ is a rational group of rank $1$. 
\end{proof}

\vskip2mm
\begin{lemma}\label{Lemma 5.33}
Let $A$ be a split local commutative graded by the radical algebra over $k$ with Lowey length $l$ and $\emph{dim}(J/J^2)=n$. Then the admissible ideal $I$ contains an $\emph{\text{Im}}(\Phi_A)$-stable finite-dimensional subspace $W$ with all non-zero elements having the same degree.
\end{lemma}
\vskip1mm
\begin{proof}
    Since $A$ is a graded by the radical local algebra, $I$ is isomorphic to $\langle X_1,\dots,X_n\rangle^l+\langle f_1,\dots,f_m\rangle$ with $f_i$' s  as homogeneous polynomials of $2\leq \text{deg}(f_i)\leq (l-1)$. Let $S$ denote the collection of polynomials of minimal degree in the generating set of $I$; \[S=\{f_i: \text{deg }f_i\leq \text{deg }f_j, \forall i\neq j \text{ and }i,j\in\{1,\dots,m\}\}.\] By renumbering if necessary, we can assume that $S=\{f_1,\dots,f_r\}$, where $r\leq m$. We claim that $W=\text{ Span}(S)$ is an $\text{Im }(\Phi_A)$-stable subspace under the action defined by $F_1.f=F_1(f)$, for all $F_1\in \text{Im}(\Phi_A)$ and $f\in k[X_1,\dots,X_n]$, where $F_1(f)(X)=f(XM_{F_1})$. We have $W=\{c_1f_1+\dots+c_rf_r:c_i\in k \}$. It is known that $F_1(I)\subset I$ for all $F_1\in \text{Im}(\Phi_A)$, by Lemma \ref{Lemma 5.3}, where $F_1$ is given by a linear change of variables. Let $0\neq g\in W$, then $F_1(g)\in I$. We have $\text{deg}(F_1(g))=\text{deg}(g)$ and $F_{1}(g)$ is a non-zero homogeneous polynomial as $F_1$ is given by a linear change of variables. Hence, $\text{deg}(g)=\text{deg}(F_{1}(g))=\text{Max}_{j=1}^m(\text{deg}(g_j)+\text{deg}(f_j))=\text{deg}(g_j)+\text{deg}(f_j)$, where $F_1(g)=\sum_{j=1}^{m}g_jf_j$, $g_j \text{ is a homogeneous polynomial in } k[X_1,\dots,X_n]$ and  $j\in\{1,\dots,m\}$. However, $\text{deg}(g)=\text{deg}(f_i)<\text{deg}(f_s)$, $\forall i\in \{1,\dots,r\},\forall s\in \{r+1,\dots,m\}$. Therefore, $\text{deg}(g_i)=0$ and $g_s=0$. This implies $F_1(g)=a_1f_1+a_2f_2+\dots+a_rf_r\in W$, where $a_i\in k.$
\end{proof}
\vskip2mm
\begin{remark}\label{Remark 5.34}
    The above subspace $W$ is independent of the choice of generators of $I$. We call this subspace the \emph{minimal degree subspace} associated to $A$ as it is generated by polynomials of minimal degree in $I$ and $I$ is determined by $A$.
\end{remark}
\vskip2mm
\begin{theorem}\label{Theorem 5.35}
Let $A$ be a split local commutative graded by the radical algebra over a field $k$ with $\emph{dim}(J/J^2)=n$. Let $W$ be the minimal degree subspace associated to $A$. Assume $W$ contains a non-singular homogeneous polynomial of degree at least $3$ if $\emph{char }k=0$. Additionally if $\emph{char }k=p>3$ and $W$ contains a non-singular homogeneous polynomial of degree $d$ with $2<d<p$. Then the following two statements are equivalent: 
\begin{enumerate}
\item $G_A$ is $k$-split solvable; 
\vskip2mm
\item $W$ has a stable full flag by $\emph{\text{Im}}(\Phi_A)$.
\end{enumerate}
Therefore, under the condition $(2)$, the group is rational. In particular, if $\text{dim }W=1$, then $G_A$ is rational.
\end{theorem}
\vskip1mm
\begin{proof}
By Lemma \ref{Lemma 5.33}, $W$ is a stable subspace under $\text{Im}(\Phi_A$). Consider the flag variety of $W$, i.e., $\mathfrak{G}(W)$, which is a complete variety with the action of $\text{Im}(\Phi_A$). However, $G_A$ is $k$-split solvable if and only if $\text{Im}(\Phi_A$) is $k$-split solvable. Using Proposition $15.2$ of \cite{Borel}, we get a fixed point in $\mathfrak{G}(W)$. Therefore, we get a stable full flag of $W$. Conversely, if $\text{Im}(\Phi_A)$ has a stable full flag in $W$, then the representation of $\text{Im}(\Phi_A)$ in $\text{GL}(W)$ is $k$-split solvable. Let $s\in W$ be a non-singular homogeneous polynomial such that $\text{deg}(s)\geq 3$ if $\text{char }k=0$, and if $\text{char }k=p>3$ then $s\in W$ be a non-singular homogeneous polynomial of degree $d$ with $2<d<p$. Let $B$ be a basis of $W$ containing $s$. Then use the basis $B$ to represent $\text{Im}(\Phi_A)$ in $\text{GL}(W)$ in the following map $\Psi$, where $\Psi(M)(g)=g(XM)$ and $g\in B$. So we have the representation \[\Psi:\text{Im}(\Phi_A)\rightarrow \text{GL}(W)\] with $\text{Ker}(\Psi)\subset \text{Stab}^{0}(s)$. However, by Theorem \ref{Theorem 5.31}, $\text{Stab}^{0}(s)=\{Id\}$. Therefore, $\text{Im}(\Phi_A)\cong \text{Im}\Psi.$ Since $\text{Im}(\Phi_A)$ has a stable full flag in $W$, $\text{Im}\Psi$ is conjugate to a $k$-split solvable group. Hence, $\text{Im}(\Phi_A)$ is $k$-split solvable.
\end{proof}

\section{Counter-Example to $R$-triviality of $G_A$}\label{counter}
\vskip2mm 
\noindent
In this section, we construct a counter-example of Question \ref{Question 5.8} using the example in the paper \cite{Gille2} and Lemma $5.1$ of \cite{GG}. The following are well-known examples of non $R$-trivial groups and theorems associated with them. 
\vskip2mm
\noindent
Merkurjev in \cite{AM2} constructs a quadratic form $q$ of dimension $6$ with non-trivial signed discriminant over a base field $k$ of characteristic not $2$ and cohomological dimension $2$ such that $\text{PSO}(q)$ is not $R$-trivial as a variety, where $\text{PSO}(q)$ is projective special orthogonal group.

\vskip2mm
\begin{theorem}\label{Theorem 6.1}
$($\cite{Gille2}$)$ There exists a field $k$ of characteristic $0$ with cohomological dimension $3$ and a quadratic form $q/k$ with rank 8 and trivial signed discriminant such that the variety $\emph{PSO}(q)$ is not stably $k$-rational.
\end{theorem}
\par 
\noindent
In addition, Ph.~Gille showed that the above group $\text{PSO}(q)$ is not $R$-trivial. Using this theorem, Nivedita Bhaskhar has produced more examples of adjoint non-$R$ trivial groups in \cite{Nivedita}.


\vskip2mm

\begin{theorem}\label{Theorem 6.2}
$($\cite{Nivedita}, \emph{Theorem} $7.1)$ For each $n=2m (m\geq 3)$, there exists a quadratic form $q_n$ defined over a field $k_n$ such that $\emph{PSO}(q_n)$ is not $R$-trivial as $k_n$ group.
\end{theorem}

\vskip2mm

\begin{theorem}\label{Theorem 6.3}
$($\cite{Preeti}, \emph{Theorem} $1.1)$ Let $p$ be a prime such that $p\neq 2$. Let $F=Q_p(t)$ be the rational function field in one variable over the $p$-adic field $Q_p$. Then for every positive integer $n\geq 2$, there exist absolutely simple adjoint algebraic groups G of type $^{2}\emph{A}_{2n-1}$
 over $F$ such that, the group of rational equivalence classes over $F$ is non-trivial,i.e., $G(F)/R \neq (1)$.
\end{theorem}
\vskip 2mm
\noindent
Now, we are in a position to construct our counter-example. We proceed as follows.

\vskip2mm
\noindent
As in Theorem \ref{Theorem 6.1}, there exists a field $k$ of characteristic $0$ with cohomological dimension equal to $3$ and a quadratic form $q$ with the given property as in Theorem \ref{Theorem 6.1} such that $\text{PSO}(q)(F)/R\neq (1)$ for some field extension $F/k$. It is now possible to apply Lemma \ref{Lemma 5.29}. Prior to this, it suffices to identify an irreducible representation of $\text{PSO}(q)$, as any irreducible representation is restricted and tensor indecomposable when the characteristic of $k$ is $0$. Let us denote such an irreducible representation by $V$, ensuring that the pair $(\text{PSO}(q), V)$ is not listed in Tables $C$ and $D$ of the article \cite{GG}. Hence, according to our choice $k[V]^{\text{PSO}(q)}\neq k$ or $k[q]$. Consequently, by applying Lemma \ref{Lemma 5.29}, we can identify a homogeneous polynomial \( f \) such that \( f \in k[V]^{\text{PSO}(q)} - k[q] \) and the identity component of the stabilizer of \( f \) is \( \text{PSO}(q) \). Now we have enough data to construct a suitable commutative associative algebra $A$ such that $G_A$  is not $R$-trivial. Consider the following algebra $A=k[V]/I\cong k[X_1,\dots,X_m]/I$, where $I=\langle X_1,\dots,X_m\rangle^l+\langle f\rangle$ and $m=\text{dim }V, l>\text{deg}(f)>2$.
\vskip2mm 
\noindent
Then we claim that $G_A$ (the identity component of $\text{Aut}_k(A)$) is not $R$-trivial.  It is clear that $A$ is graded by the radical algebra. Therefore, Proposition \ref{Proposition 5.4} gives  
\[\text{Aut}_k(A)\cong \text{Im}(\Phi_A) \ltimes \text{U}, \quad \text{ where } U\text{ is a connected unipotent group.}\]
The constructed algebra $A$ satisfies the hypothesis of Proposition \ref{Proposition 5.12}, hence $\text{Im}(\Phi_A)=\text{Sim}(f)=\mathbb{G}_m.\text{Stab}(f)$. Therefore,
$G_A\cong \mathbb{G}_m.\text{Stab}(f)^{0}\ltimes \text{U}.$
As $\text{Stab}(f)^0=\text{PSO}(q)$, we get $\mathbb{G}_m.\text{Stab}(f)^{0}\cong(\mathbb{G}_m\times \text{PSO}(q))/\mathbb{G}_m\cap \text{PSO}(q)$. It is clear that $\mathbb{G}_m\cap \text{PSO}(q)=Z(\text{PSO}(q))=\{Id\}$. So we have 
\begin{equation*}
G_A(F)/R=\text{PSO}(q)(F)/R\neq (1) \text{ for some } F/k.
\end{equation*}
\vskip2mm
\noindent
 Following the same technique with Merkurjev's example, where $q$ is of dimension $6$ and has non-trivial signed discriminant over a field $k$ with characteristic $k=0$ and cohomological dimension $2$, it is possible to construct an algebra $A$ such that the dimension of $J/J^2$ is at least $6$, and the group $G_A$ is not $R$-trivial.   
The application of Theorem \ref{Theorem 6.2} and Theorem \ref{Theorem 6.3}, along with comparable techniques discussed earlier, facilitates the construction of additional examples of non $R$-trivial groups \( G_A \). In particular, for all integers \( i \) such that \( i\geq 4 \), Theorem \ref{Theorem 6.2} can be utilized to establish a non $R$-trivial group \( G_A \) with a rank equal to \( i \) and $\text{dim }(J/J^2)\geq 6$. As a consequence of the preceding discussion, we have the following theorem.
\begin{theorem}\label{Theorem 6.4}
For each $n\geq3$, there exists a finite-dimensional commutative algebra $A_n$ with $\emph{dim } \displaystyle (\frac{J_n}{J_n^2})\geq 2n$, where $J_n$ is the Jacobson radical of $A_n$ over a field $k_n$, such that $G_{A_n}$ is not $R$-trivial as a $k_n$-group and rank of $G_{A_n}$ is $(n+1)$. 
\end{theorem}
\vskip1mm
\begin{proof}
Let $n\geq 3$. By Theorem \ref{Theorem 6.2} for each $2n$ there exists a quadratic form $q_{2n}$ defined over a field $k_n$ such that $\mathrm{PSO}(q_{2n})$ is not $R$-trivial as $k_n$-group. Now applying Lemma \ref{Lemma 5.29} to each pair $(\mathrm{PSO}(q_{2n}),V_n)$ over the field $k_n$ with $k_n[V_n]^{\mathrm{PSO}(q_{2n})}\neq k \text{ or }k[q_{2n}]$, we get that there exists a homogeneous polynomial $f_n\in k[V_n]^{\mathrm{PSO}(q_{2n})}$ for each $n\geq 3$ such that $\mathrm{PSO}(q_{2n})=\mathrm{Stab}^{0}(f_n)$. Hence, similarly as above we consider the algebra $A_n=k[V_n]/I_n$, where $I_n=\langle X_1,\dots,X_{m_{n}}\rangle ^{l_n}+\langle f_n\rangle$ and $\text{ dim }V_n=m_n\geq 2n$, $l_n> \text{deg}(f_n)>2$. Therefore, we have \[G_{A_n}(k_n)/R= \mathrm{PSO}(q_{2n})(k_n)/R\neq (1).\]
\end{proof}
\vspace{2mm}
\noindent
We have only given some sufficient conditions on $A$ so that $G_A$ is $R$-trivial for a split finite-dimensional associative algebra $A$. Still, the following question remains open.
\begin{question}
Let $A$ be a split finite-dimensional associative algebra over a perfect field $k$. What is a necessary and sufficient condition on the algebra $A$, so that $G_A$ is $R$-trivial or rational? 
\end{question}
 \noindent
 {\bf Acknowledgments:} I would like to thank my supervisor, Prof. Maneesh Thakur, for suggesting this problem. I am also grateful for his guidance and fruitful discussions that I had with him during the preparation of this manuscript. I thank Sayan Pal and Arghya Chongdar for their useful suggestions. I also thank Indian Statistical Institute for financial support during the work.



\bibliography{references}
\bibliographystyle{amsplain}
\end{document}


